\exercisesfor{5}

第 5 节的约定除非另有说明均有效。

\begin{enumerate}
\def\labelenumi{\arabic{enumi}.}
\item
  设 g 为 2 维非交换可解李代数。证明 g 的 Killing 形式非零，g 上的每个不变双线性形式都退化，且 g 的每个导子都是内导子。
\item
  \begin{enumerate}
  \def\labelenumii{(\alph{enumii})}
  \tightlist
  \item
    证明在 \(\mathbf{R}\) 上由乘法表 \([x, y] = z\)、\([x, z] = -y\)、\([y, z] = 0\) 定义的 3 维可解李代数中，不存在维数为 3、2、1、0 的递降理想序列。
  \end{enumerate}
\end{enumerate}

\begin{enumerate}
\def\labelenumi{(\alph{enumi})}
\setcounter{enumi}{1}
\tightlist
\item
  证明在非交换的 2 维可解李代数 g 中存在维数为 2、1、0 的理想序列，但 g 并非幂零。
\end{enumerate}

\begin{enumerate}
\def\labelenumi{\arabic{enumi}.}
\setcounter{enumi}{2}
\item
  设 \(\mathfrak{g}\) 为可解李代数，使得条件 \(x\in \mathfrak{g},y\in \mathfrak{g}\)、\([[x,y],y] = 0\) 蕴含 \([x,y] = 0\)。证明 \(\mathfrak{g}\) 是交换的。（设 \(k\) 为使 \(\mathcal{D}^{k - 1}\mathfrak{g}\neq \{0\}\)、\(\mathcal{D}^k\mathfrak{g} = \{0\}\) 的最大整数。假设 \(k\geqslant 2\)，先证明 \([\mathcal{D}^{k - 2}\mathfrak{g},\mathcal{D}^{k - 1}\mathfrak{g}] = \{0\}\)，再证明 \([\mathcal{D}^{k - 2}\mathfrak{g},\mathcal{D}^{k - 2}\mathfrak{g}] = \{0\}\)，从而得出矛盾。）
\item
  证明 \(\mathfrak{st}(n,\mathbf{K})\) 的中心为零，而 \(\mathfrak{n}(n,\mathbf{K})\) 的中心维数为 1。
\item
  设 \(\mathfrak{g}\) 为李代数，\((\mathcal{D}^0\mathfrak{g},\mathcal{D}^1\mathfrak{g},\ldots ,\mathcal{D}^n\mathfrak{g})\) 为 \(\mathfrak{g}\) 的导代数序列（ \(n\geqslant 0,\mathcal{D}^{n - 1}\mathfrak{g}\neq \mathcal{D}^n\mathfrak{g}\) ）。则 \(\dim \mathcal{D}^i\mathfrak{g} / \mathcal{D}^{i + 1}\mathfrak{g}\geqslant 2^{i - 1} + 1\) 对 \(1\leqslant i\leqslant n - 2\) 成立。（对 \(\mathcal{D}^n\mathfrak{g}\) 取商，归结为 \(\mathfrak{g}\) 可解的情形。然后利用 \(\mathcal{D}_{\mathfrak{g}}\) 幂零这一事实以及第 4 节练习 7 (c)。）
\item
  \begin{enumerate}
  \def\labelenumii{(\alph{enumii})}
  \tightlist
  \item
    验证以下乘法表定义了一个 5 维可解李代数 g：
  \end{enumerate}
\end{enumerate}

\[
\begin{array}{r} \left[ x _ {1}, x _ {2} \right] = x _ {5}, \qquad \left[ x _ {1}, x _ {3} \right] = x _ {3}, \qquad \left[ x _ {2}, x _ {4} \right] = x _ {4} \\ \left[ x _ {1}, x _ {4} \right] = \left[ x _ {2}, x _ {3} \right] = \left[ x _ {3}, x _ {4} \right] = \left[ x _ {5}, \mathfrak {g} \right] = 0. \end{array}
\]

\begin{enumerate}
\def\labelenumi{(\alph{enumi})}
\setcounter{enumi}{1}
\item
  证明 \(\mathfrak{g}\) 关于 Killing 形式的正交子空间是 \(\mathcal{D}\mathfrak{g} = \mathrm{K}x_3 + \mathrm{K}x_4 + \mathrm{K}x_5\)。由此推出 \(\mathcal{D}\mathfrak{g}\) 是 \(\mathfrak{g}\) 的最大幂零理想。
\item
  证明不存在 \(\mathcal{D}\mathfrak{g}\) 在 \(\mathfrak{g}\) 中的补子代数。由此推出 \(\mathfrak{g}\) 不是交换代数与一个幂零理想的半直积。（证明该幂零理想必然就是 \(\mathcal{D}\mathfrak{g}\)。）
\end{enumerate}

\begin{enumerate}
\def\labelenumi{\arabic{enumi}.}
\setcounter{enumi}{6}
\item
  设 \(\mathfrak{g}\) 为以 \((x,y,z)\) 为基、满足 \([x,y] = z\)、\([x,z] = y\)、\([y,z] = 0\) 的 3 维可解李代数。证明将 \(x\) 映为 \(-x, y\) 映为 \(-z, z\) 映为 \(y\) 的线性映射是 \(\mathfrak{g}\) 的一个 4 阶自同构。将此结果与第 4 节练习 21 (c) 相比较。
\item
  \begin{enumerate}
  \def\labelenumii{(\alph{enumii})}
  \tightlist
  \item
    设 \(\mathfrak{g}_0\) 为 3 维实可解李代数，使得 \(\mathcal{D}\mathfrak{g}_0\) 交换且 2 维。设 \(\mathfrak{g}\) 为由 \(\mathfrak{g}_0\) 通过将基域从 \(\mathbf{R}\) 扩张到 \(\mathbf{C}\) 而得到的代数。对 \(x \in \mathfrak{g}\)，令 \(u_x\) 为 \(\mathrm{ad}_{\mathfrak{g}}x\) 在 \(\mathcal{D}\mathfrak{g}\) 上的限制。证明 \(u_x\) 的特征值或者具有相同的绝对值，或者在 \(\mathbf{R}\) 上线性相关。（我们有 \(x = \lambda y + z\)，其中 \(z \in \mathcal{D}\mathfrak{g}\)、\(y \in \mathfrak{g}_0\)、\(\lambda \in \mathbf{C}\)，因而 \(u_x = \lambda u_y\)。而 \(u_y\) 是 \(\mathbf{C}\) 线性地扩张 \(\mathcal{D}\mathfrak{g}\) 的一个 \(\mathbf{R}\) 线性自同态所得的映射，后者定义在 \(\mathcal{D}\mathfrak{g}_{0}\) 上。）
  \end{enumerate}
\end{enumerate}

\begin{enumerate}
\def\labelenumi{(\alph{enumi})}
\setcounter{enumi}{1}
\item
  证明存在一个 3 维复可解李代数 g，其 \(\mathcal{D}_{\mathfrak{g}}\) 交换且 2 维，并且存在 g 的元素 \(x\)，使得 \(\mathrm{ad}_{\mathfrak{g}}x\) 在 \(\mathcal{D}_{\mathfrak{g}}\) 上的限制的特征值既不具有相同的绝对值，也不在 R 上线性相关。（将 g 构造为 1 维代数与 2 维交换代数的半直积。）
\item
  证明 (b) 中构造的代数不能由某个实李代数通过将基域从 \(\mathbf{R}\) 扩张到 \(\mathbf{C}\) 而得到。
\end{enumerate}

\begin{enumerate}
\def\labelenumi{\arabic{enumi}.}
\setcounter{enumi}{8}
\tightlist
\item
  设 \(\mathfrak{g}\) 为李代数，\(\mathfrak{r}\) 为其根基，\(\mathfrak{n}\) 为其最大幂零理想，D 为 \(\mathfrak{g}\) 的一个导子。证明 \(\mathcal{D}\mathfrak{g} \cap \mathfrak{r} \subset \mathfrak{n}\)。（设 \(\mathfrak{d}\) 为 \(\mathfrak{g}\) 的导子组成的李代数，\(\mathfrak{r}'\) 为其根基。若 \(x \in \mathfrak{g}\) 使得 \(\mathrm{D}x \in \mathfrak{r}\)，则
\end{enumerate}

\[
\operatorname{ad} (\mathrm{D} x) = [ \mathrm{D}, \operatorname{ad} x ] \in \mathscr {D} \mathfrak {d} \cap \mathfrak {r} ^ {\prime}
\]

（命题 5 的推论 2），因而 \(\operatorname{ad}(\mathbf{D}x)\) 幂零（定理 1），从而 \(\mathbf{D}x \in \mathfrak{n}\)。）

\begin{enumerate}
\def\labelenumi{\arabic{enumi}.}
\setcounter{enumi}{9}
\item
  设 g 为李代数，r 为其根基，a 为 g 的一个次不变子代数。证明 a 的根基是 \(a \cap r\)。（若干次应用命题 5 的推论 3。）
\item
  设 \(\mathfrak{g}\) 为李代数，\(\rho\) 为 \(\mathfrak{g}\) 的有限维表示，\(\mathbb{E}\) 为由 1 与 \(\rho(\mathfrak{g})\) 生成的结合自同态代数。证明理想 \(n\)（\(\rho\) 的最大幂零性理想）等于集合 \(n'\)，它由 \(x \in \mathfrak{g}\)（满足 \(\operatorname{Tr}(\rho(x)u) = 0\)，对一切 \(u \in E\)）组成。（为证明 \(n' \subset n\)，证明 \(n'\) 是一个理想，且对一切 \(x \in n'\)，\(\rho(x)\) 的半单分量为零，这通过注意 \(\operatorname{Tr}((\rho(x))^n) = 0\) 对每个整数 \(n > 0\) 成立得到。）
\item
  假设 K 代数闭。设 g 为幂零李 K-代数。设 \(\rho\) 为 g 在向量空间 V 上的一个有限维表示。对 g 上的每个线性形式 \(\lambda\)，令 \(V^{\lambda}\) 为 V 中由满足如下条件的 \(\xi \in V\) 组成的向量子空间：对一切 \(x \in g\)，当 n 充分大时 \((\rho(x) - \lambda(x)\mathrm{I})^{n}\xi = 0\)。
\end{enumerate}

\begin{enumerate}
\def\labelenumi{(\alph{enumi})}
\item
  诸 \(V^{\lambda}\) 在 \(\rho(\mathfrak{g})\) 下稳定，且诸 \(V^{\lambda}\) 之和是直和。（使用第 4 节练习 22。）
\item
  我们有 \(\mathrm{V} = \sum \mathrm{V}^{\lambda}\)。（若每个 \(\rho(x)\) 都只有一个特征值，则由定理 1 的推论 2 得 \(\mathrm{V} = \mathrm{V}^{\lambda_0}\)。若 \(\rho(x_0)\) 至少有两个不同的特征值，则 \(\mathrm{V}\)）
\end{enumerate}

是两个在 \(\rho(\mathfrak{g})\) 下稳定的非平凡子空间的直和。然后对 V 的维数作归纳。）

\begin{enumerate}
\def\labelenumi{\arabic{enumi}.}
\setcounter{enumi}{12}
\item
  设 \(\mathfrak{g}\) 为李代数，\(\mathfrak{D}\) 为 \(\mathfrak{g}\) 的导子组成的李代数。\(\mathfrak{g}\) 特征幂零的充分必要条件是 \(\mathfrak{D}\) 幂零且 \(\dim \mathfrak{g} > 1\)。（为验证条件充分，把 \(\mathfrak{g}\) 写成子空间 \(\mathfrak{g}^{\lambda}\) 之和，办法是把练习 12 应用于 \(\mathfrak{D}\) 的恒等映射。证明 \([\mathfrak{g}^{\lambda}, \mathfrak{g}^{\mu}] \subset \mathfrak{g}^{\lambda + \mu}\)，且每个 \(\mathfrak{g}^{\lambda}\) 都是 \(\mathfrak{g}\) 的理想。由此推出 \(\mathfrak{g}^{\lambda}\) 当 \(\lambda \neq 0\) 时是交换的。再利用 \(\mathfrak{D}\) 幂零这一事实，证明 \(\mathfrak{g} = \mathfrak{g}^{0}\)（若 \(\dim \mathfrak{g} > 1\)）；否则 \(\mathfrak{g} = \mathfrak{g}^{0} \times \mathfrak{h}\)，其中 \(\mathfrak{h}\) 交换。先证明 \(\dim \mathfrak{h} \leqslant 1\)。若 \(\dim \mathfrak{h} = 1\)，注意将存在 \(\mathfrak{g}\) 的导子 D，使得 D( \(\mathfrak{g}^{0}\) ) = \{0\} 且 D( \(\mathfrak{h}\) ) 包含在 \(\mathfrak{g}^{0}\) 的中心中（ \(\S 4\)，练习 8 ( \(a\) )）。）
\item
  设 V 为 K 上有限维向量空间，z 为 V 的自同态。我们采用第 3 节练习 3 的记号 \(z_{q}^{p}\)。V 的自同态 \(z'\) 若对一切 p 和 q 使得 \(z_{q}^{p}\) 的每个零点都是 \(z_{q}^{\prime p}\) 的零点，则称为 z 的一个复本。证明若 \(\operatorname{Tr}(zz') = 0\) 对每个复本 \(z'\)（\(z\) 的复本）成立，则 z 幂零。（利用引理 3 的证明。在该证明的记号下，特别证明 t 是 z 的复本。）
\item
  设 K 为特征 2 的域。幂零李代数 \(\mathfrak{sl}(2,\mathbf{K})\) 在 \(K^{2}\) 上的恒等表示定义了 \(\mathfrak{sl}(2,\mathbf{K})\) 与 \(K^{2}\) 的一个半直积 h。证明 h 可解，但 Dh 不幂零。由此推出 h 不容许任何以三角矩阵给出的忠实线性表示。再证明练习 5 的结论在此为假。
\item
  设 g 为任意域 K 上的李代数，A 为 K 上交换结合代数，而 \(g' = g \otimes_{K} A\)，它可视为 K 上的李代数。
\end{enumerate}

\begin{enumerate}
\def\labelenumi{(\alph{enumi})}
\item
  若 D 是 A 的导子，证明存在唯一的导子 \(D'\)（\(g'\) 的导子），使得 \(\mathrm{D}'(x \otimes a) = x \otimes \mathrm{D}a\) 对 \(x \in g, a \in A\) 成立。
\item
  设 \(p\) 为素数，\(G\) 为 \(p\) 阶循环群，\(s\) 为 \(G\) 的一个生成元。假设 \(K\) 的特征为 \(p\)，以下设 \(A\) 为 \(G\) 在 \(K\) 上的代数。证明存在唯一的导子 \(D\)（\(A\) 的导子），使得 \(D(s^k) = ks^{k-1}\) 对一切 \(k \in \mathbf{Z}\) 成立。证明取 \(K\) 中系数时诸 \(x \otimes (s - 1)^k\)（ \(k = 1, 2, \ldots, p - 1, x \in \mathfrak{g}\) ）的线性组合构成一个可解理想 \(\tau\)（\(g'\) 的理想），且 \(g'/\tau\) 同构于 \(\mathfrak{g}\)。
\item
  取 \(\mathfrak{g}\) 为单代数（参见~第 6 节，第 2 号，定义 2）。则 \(\mathfrak{r}\) 是 \(\mathfrak{g}'\) 的根基，但不是特征理想。（注意 \(\mathrm{D}(x\otimes (s - 1)) = 1\otimes x.\) ）
\end{enumerate}

\begin{enumerate}
\def\labelenumi{\arabic{enumi}.}
\setcounter{enumi}{16}
\tightlist
\item
  设 g 为李代数。假设对每个 K 上有限维单 g-模 M，诸 \(x_{M}\) 两两可交换。证明 g 可解。（注意 Dg 包含在幂零根基中，因而可解。）
\end{enumerate}

\exercisesfor{6}

第 6 节的约定除非另有说明均有效。

¶ 1. 设 g 为半单李代数，ρ 为 g 在 M 上的一个有限维表示。

\begin{enumerate}
\def\labelenumi{(\alph{enumi})}
\item
  若 \(\rho\) 单且非零，则 \(\mathrm{H}^p (\mathfrak{g},\mathbf{M}) = \{0\}\) 对一切 \(p\) 成立。（使用第 3 节命题 1 和练习 12 (j)。）
\item
  对一切 \(\rho\)，有 \(\mathbf{H}^1(\mathfrak{g}, \mathbf{M}) = \{0\}\)。（若 \(\rho\) 单且非零，应用 (a)。若 \(\rho\) 为零，利用 \(\mathfrak{g} = \mathcal{D}\mathfrak{g}\) 这一事实。一般情形，对 \(\rho\) 的维数作归纳：若 N 是 M 的异于 \(\{0\}\) 与 M 的子 g-模，利用正合列：
\end{enumerate}

\[
\mathrm{H} ^ {1} (\mathfrak {g}, \mathrm{N}) \rightarrow \mathrm{H} ^ {1} (\mathfrak {g}, \mathrm{M}) \rightarrow \mathrm{H} ^ {1} (\mathfrak {g}, \mathrm{M} / \mathrm{N})
\]

（此正合列已在第 3 节练习 12 (f) 中建立。）由此重新得出第 2 号的注 2。

\begin{enumerate}
\def\labelenumi{(\alph{enumi})}
\setcounter{enumi}{2}
\item
  由 (b) 推出定理 2 的一个证明。（使用第 3 节练习 12 (g)。）再由 (b) 推出 g 的每个导子都是内导子。（使用第 3 节练习 12 (h)。）
\item
  对一切 \(\rho\)，有 \(\mathrm{H}^2(\mathfrak{g}, \mathrm{M}) = \{0\}\)。（仿照 (b) 的论证，只需考虑 \(\rho = 0\) 的情形。设 \(c \in \mathrm{H}^2(\mathfrak{g}, \mathrm{M})\)。按第 3 节练习 12 (i)，考虑中心扩张 \(\mathfrak{h}\)，即 \(\mathfrak{g}\) 被 \(\mathrm{M}\) 的、由 \(c\) 定义的扩张。\(\mathfrak{h}\) 的伴随表示定义了 \(\mathfrak{g}\) 在 \(\mathfrak{h}\) 上的一个表示。由定理 2，\(\mathrm{M}\) 在 \(\mathfrak{h}\) 中有一个在 \(\mathfrak{g}\) 下稳定的补，因而该扩张是平凡的。于是由第 3 节练习 12 (i) 得 \(c = 0\)。）
\item
  由 (b) 与 (d) 推出定理 5 的一个证明。（与正文中一样，可归结为根基交换的情形。）
\end{enumerate}

\begin{enumerate}
\def\labelenumi{\arabic{enumi}.}
\setcounter{enumi}{1}
\item
  设 g 为李代数，r 为其根基，\((a_{0}, a_{1}, \ldots)\) 为 g 的理想序列，定义如下：(1) \(a_{0} = \{0\}\)；(2) \(a_{i+1}/a_{i}\) 是 \(g/a_{i}\) 的一个极大交换理想。设 p 为使 \(a_{p} = a_{p+1} = \cdots\) 的最小整数。证明 \(r = a_{p}\)。（证明 \(g/a_{p}\) 半单。）
\item
  李代数 g 半单的充分必要条件是：在每个包含 g 作为子代数的李代数中 g 都是约化的。（若 g 满足此条件，设 M 为 K 上有限维 g-模。把 M 视为交换代数，作 g 与 M 的半直积，其中 g 是约化的。由此推出 M 半单。）
\item
  设 g 为半单李代数，\(\rho\) 为 g 在有限维空间 M 上的非零单表示。设 h 为相应的半直积。证明 \(h = Dh\)，h 的中心为零，且 h 不是半单代数与可解代数的乘积。
\item
  设 \(\mathfrak{a}\) 为李代数，\(\mathfrak{r}\) 为其根基。若 \(\mathfrak{r}\) 具有特征理想的递降序列 \(\mathfrak{r} = \mathfrak{r}_0 \supset \mathfrak{r}_1 \supset \dots \supset \mathfrak{r}_n = \{0\}\)，使得 \(\dim \mathfrak{r}_i / \mathfrak{r}_{i+1} = 1\) 对 \(0 \leqslant i < n\) 成立，则 \(\mathfrak{g}\) 是半单代数与可解代数的乘积。
\end{enumerate}

（设 s 为 g 的一个 Levi 子代数。对一切 \(x \in s\)，令 \(\rho(x)\) 为 \(ad_{g}x\) 在 r 上的限制。则 \(\rho\) 是一维表示的直和，这些表示都是零表示，因为 \(s = Ds\)。）

\begin{enumerate}
\def\labelenumi{\arabic{enumi}.}
\setcounter{enumi}{5}
\item
  设 \(\mathfrak{g}\) 为李代数，\(\mathcal{D}^{\infty}\mathfrak{g}\) 为诸 \(\mathcal{D}^{p}\mathfrak{g}\)（ \(p = 1,2,\ldots\) ）的交。证明 \(\mathfrak{g}/\mathcal{D}^{\infty}\mathfrak{g}\) 可解，且 \(\mathcal{D}(\mathcal{D}^{\infty}\mathfrak{g}) = \mathcal{D}^{\infty}\mathfrak{g}\)。\(\mathfrak{g}\) 同构于半单代数与可解代数的乘积的充分必要条件是 \(\mathcal{D}^{\infty}\mathfrak{g}\) 半单。
\item
  \begin{enumerate}
  \def\labelenumii{(\alph{enumii})}
  \tightlist
  \item
    设 g 为李代数，b 为 g 的半单理想，a 为 b 在 g 中的中心化子，于是 g 与 \(b \times a\) 恒同。证明对 g 的每个理想 \(\ell\)，有 \(\ell = (\ell \cap \mathfrak{b}) \times (\ell \cap \mathfrak{a})\)。（设 \(\ell_{1}\) 为 \(\ell\) 到 b 上的典范投影；它是 b 的理想，因而 \(\mathcal{D}\ell_{1} = \ell_{1}\)；由此推出 \(\ell_{1} \subset \ell\)。）
  \end{enumerate}
\end{enumerate}

\begin{enumerate}
\def\labelenumi{(\alph{enumi})}
\setcounter{enumi}{1}
\item
  设 \(\beta\) 为 \(\mathfrak{g}\) 上的不变双线性形式。证明 \(\mathfrak{b}\) 与 \(\mathfrak{a}\) 关于 \(\beta\) 正交（利用 \(\mathfrak{b} = \mathcal{D}\mathfrak{b}\) 以及 \(\beta = \beta_{1} + \beta_{2}\) 这一事实，其中 \(\beta_{1}\)（相应地 \(\beta_{2}\) ）是到 \(\mathfrak{a}\)（相应地 \(\mathfrak{b}\) ）上的限制为零的不变双线性形式。
\item
  由 (a) 推出 \(\mathfrak{g}\) 中存在一个最大半单理想。（考虑 \(\mathfrak{g}\) 的一个极大半单理想。）
\end{enumerate}

\begin{enumerate}
\def\labelenumi{\arabic{enumi}.}
\setcounter{enumi}{7}
\tightlist
\item
  设 g 为李代数。若 \(h \neq \{0\}\) 且 g 的包含于 h 中的每个理想都等于 \(\{0\}\) 或 h，则称 g 的理想 h 为极小理想。
\end{enumerate}

\begin{enumerate}
\def\labelenumi{(\alph{enumi})}
\item
  g 的每个单理想都是极小的。
\item
  设 \(\mathfrak{h}\) 为 \(\mathfrak{g}\) 的极小理想，\(\mathfrak{r}\) 为 \(\mathfrak{g}\) 的根基。则或者 \(\mathfrak{h} \subset \mathfrak{r}\)，此时 \(\mathfrak{h}\) 交换；或者 \(\mathfrak{h} \cap \mathfrak{r} = \{0\}\)，此时 \(\mathfrak{h}\) 单。（利用李代数的导理想与半单李代数的单分量都是特征理想这一事实。）
\end{enumerate}

\begin{enumerate}
\def\labelenumi{\arabic{enumi}.}
\setcounter{enumi}{8}
\item
  李代数 g 约化的充分必要条件是其中心 c 等于它的最大幂零理想。（若该条件成立，设 r 为 g 的根基；Dr 包含在 r 的中心中，因而 r 幂零，从而 r = c。）
\item
  设 \(\mathfrak{g}\) 为李代数，使得条件 \(x \in \mathfrak{g}, y \in \mathfrak{g}, [[x, y], y] = 0\) 蕴含 \([x, y] = 0\)。证明 \(\mathfrak{g}\) 约化。（利用第 5 节练习 3 证明根基 \(\mathfrak{r}\)（\(\mathfrak{g}\) 的根基）交换。然后证明 \([\mathfrak{g}, \mathfrak{r}] = 0\)。）
\end{enumerate}

¶ 11. (a) 设 g 为李代数，r 为其根基，s 为 g 的 Levi 子代数，m 为包含 r 的 g 的理想。则存在 s 的理想 t，在 g 中是 m 的补，使得 \([m, t] \subset r\)。

\begin{enumerate}
\def\labelenumi{(\alph{enumi})}
\setcounter{enumi}{1}
\tightlist
\item
  设 g 为李代数，a 为 g 的次不变子代数。则存在合成列 \(g = g_{0} \supset g_{1} \supset \cdots \supset g_{k} = a\)，使得 \(g_{i}\) 是 \(g_{i+1}\) 与一个子代数 \(b_{i}\) 的直和，后者或者是 1 维且包含在 \(r(g_{i})\)（即 \(g_{i}\) 的根基）中，或者是单的，且满足 \([b_{i}, g_{i+1}] \subset r(g_{i+1})\)。（归结为 a 是 g 的理想的情形。设 \(g' = g/a\)，\(r'\) 为 \(g'\) 的根基。代数 \(g'/r'\) 是它的单理想 \(a_{1}, a_{2}, \ldots, a_{c}\) 的乘积。
\end{enumerate}

取 \(\mathfrak{g}'\) 的一个合成列，它由 \(\mathbf{r}'\) 的逐次商均为 1 维的合成列以及理想 \(a_1, a_1 \times a_2, \ldots, a_1 \times a_2 \times \cdots \times a_c\) 的原像组成。然后取该合成列在 \(\mathfrak{g}\) 中的原像。）

\begin{enumerate}
\def\labelenumi{\arabic{enumi}.}
\setcounter{enumi}{11}
\tightlist
\item
  设 g 为李代数，r 为其根基，D 为 g 的一个导子。
\end{enumerate}

\begin{enumerate}
\def\labelenumi{(\alph{enumi})}
\tightlist
\item
  若 \(\mathrm{D}(\mathfrak{r}) = \{0\} \)，则 D 是内导子。（设 c 为 r 在 g 中的中心化子。g 的伴随表示定义了 \(x^{*} \mapsto \rho(x^{*})\)，即 \(g^{*} = g/r\) 在 c 上的一个表示。另一方面，\([\mathrm{D}(\mathfrak{g}), \mathfrak{r}] \subset \mathrm{D}([\mathfrak{g}, \mathfrak{r}]) + [\mathfrak{g}, \mathrm{D}(\mathfrak{r})] = \{0\}\)，因而 \(\mathrm{D}(\mathfrak{g}) \subset c\)，于是 D 定义了一个线性映射 \(D^{*}: g^{*} \to c\)。证明
\end{enumerate}

\[
\mathrm{D} ^ {*} ([ x ^ {*}, y ^ {*} ]) = \rho (x ^ {*}) \mathrm{D} ^ {*} y ^ {*} - \rho (y ^ {*}) \mathrm{D} ^ {*} x ^ {*}
\]

对一切 \(x^{*}, y^{*}\) ∈ \(\mathfrak{g}^{*}\) 成立。由此推出存在 \(a \in \mathfrak{c}\)，使得 \(D^{*}x^{*} = \rho(x^{*})a\) 对一切 \(x^{*} \in \mathfrak{g}^{*}\) 成立（参见~第 2 号，注 2），从而 \(Dx = [x, a]\) 对一切 \(x \in \mathfrak{g}\) 成立。）

\begin{enumerate}
\def\labelenumi{(\alph{enumi})}
\setcounter{enumi}{1}
\tightlist
\item
  若 D 在 r 上与 g 的某个内导子重合，则 D 是内导子。（应用 (a)。）
\end{enumerate}

\begin{enumerate}
\def\labelenumi{\arabic{enumi}.}
\setcounter{enumi}{12}
\item
  设 g 为代数闭域上的李代数，r 为其根基，\(\rho\) 为 g 的有限维单表示。则 \(\rho(z)\) 对 \(z \in r\) 是标量。（归结为 g 约化的情形。此时 r 是 g 的中心。）
\item
  设 X 为 \(\mathfrak{gl}(n, \mathbf{K})\) 中的幂零矩阵。证明在 \(\mathfrak{gl}(n, \mathbf{K})\) 中存在另外两个矩阵 Y、H，使得 \([H, X] = X\)、\([H, Y] = -Y\)、\([X, Y] = H\)。（对 n 作归纳，利用 X 的 Jordan 标准形；于是归结为 \(X = E_{12} + E_{23} + \cdots + E_{n-1,n}\) 的情形；然后取 Y 为诸 \(E_{k,k-1}\) 的线性组合，H 为诸 \(E_{jj'}\) 的线性组合。）
\end{enumerate}

¶ 15. (a) 设 X、Y 为 \(\mathfrak{gl}(n, \mathrm{K})\) 的两个矩阵。证明若 \([X, Y] = X\)，则 X 幂零。（对每个多项式 \(f \in K[T]\)，注意

\[
[ f (X), Y ] = X f ^ {\prime} (X).)
\]

\begin{enumerate}
\def\labelenumi{(\alph{enumi})}
\setcounter{enumi}{1}
\item
  设 \(\mathfrak{g}\) 为李代数，\(h\)、\(x\) 为 \(\mathfrak{g}\) 的两个元素，满足 \([h, x] = x\)，且存在 \(z \in \mathfrak{g}\) 使 \([x, z] = h\)。证明存在 \(y \in \mathfrak{g}\)，使 \([x, y] = h\) 且 \([h, y] = -y\)。（首先注意 \([z, h] + z\) 属于 \(\mathfrak{n}\)，即 \(x\) 在 \(\mathfrak{g}\) 中的中心化子。然后证明 \(\mathfrak{n}\) 在 ad \(h\) 下稳定，且 \(\mathfrak{n}\) 上的 I - ad \(h\) 的限制是双射；为此，注意到 ad \(x\) 幂零（利用 \((a)\)），并证明：若 \(\mathfrak{g}_{\mathfrak{l}}\) 是 \(\mathfrak{g}\) 在 \((\mathrm{ad}x)^{\mathfrak{l}}\) 下的像，则 I - ad \(h\) 在取商后给出 \((\mathfrak{n} \cap \mathfrak{g}_{\mathfrak{l}-1}) / (\mathfrak{n} \cap \mathfrak{g}_{\mathfrak{l}})\) 的一个双射；为此利用关系式 \([\mathrm{ad}z, (\mathrm{ad}x)^{\mathfrak{k}}] = -k(\mathrm{ad}x)^{\mathfrak{k}-1}(\mathrm{ad}h + \frac{k-1}{2}I).)\)）
\item
  设 \(\mathfrak{g}\) 为 \(\mathfrak{gl}(n, \mathbf{K})\) 的子代数；假设对每个幂零矩阵 \(X \in \mathfrak{g}\)，都存在另外两个矩阵 \(H, Y\)（均在 \(\mathfrak{g}\) 中），使得 \([H, X] = X\)、\([H, Y] = -Y\)、\([X, Y] = H\)。设 \(\mathfrak{h}\) 为 \(\mathfrak{g}\) 的子代数，使得存在
\end{enumerate}

h 在 g 中的一个补子空间 m，且满足 \([h, m] \subset m\)。证明 h 对其所有幂零矩阵具有与 g 相同的性质（利用 (b)）。

\begin{enumerate}
\def\labelenumi{\arabic{enumi}.}
\setcounter{enumi}{15}
\item
  设 g 为 \(\mathfrak{gl}(n,\mathbf{K})\) 的子代数，其恒等表示半单。证明每个幂零矩阵 \(X\in g\) 都包含在 g 的一个 3 维单李子代数中。（证明 g 在 \(\mathfrak{gl}(n,\mathbf{K})\) 中约化；然后利用练习 14 与 15 (c)。）
\item
  证明由实单李代数通过将基域从 \(\mathbf{R}\) 扩张到 \(\mathbf{C}\) 而得到的代数并不总是单的。（设 \(\mathfrak{g}\) 为实单李代数。若 \(\mathfrak{g}' = \mathfrak{g}_{(\mathbf{C})}\) 单，设 \(\mathfrak{h}\) 为由 \(\mathfrak{g}'\) 通过把标量域限制到 \(\mathbf{R}\) 而得到的实李代数。我们知道 \(\mathfrak{h}\) 是单的。于是由第 1 节练习 4，\(\mathfrak{h}_{(\mathbf{C})}\) 是两个同构于 \(\mathfrak{g}'\) 的代数的乘积。另参见练习 26 (b)。
\item
  \begin{enumerate}
  \def\labelenumii{(\alph{enumii})}
  \tightlist
  \item
    设 g 为单李代数。g 上的每个不变双线性形式或为零或非退化。若 K 代数闭，则 g 上的每个不变双线性形式 \(\beta\) 都与 Killing 形式 \(\beta_{0}\) 成比例。（考虑向量空间 g 的自同态 \(\sigma\)，它由 \(\beta(x,y)=\beta_{0}(\sigma(x),y)\) 定义；证明 \(\sigma\) 与 \(ad_{g}x\) 交换，从而是标量。）证明当 K 不是代数闭时该结论可能不成立。（利用练习 17 以及不变双线性形式空间的维数在基域扩张下不变这一事实。）
  \end{enumerate}
\end{enumerate}

\begin{enumerate}
\def\labelenumi{(\alph{enumi})}
\setcounter{enumi}{1}
\item
  设 g 为代数闭域上的半单李代数。由 (a) 与练习 7 (b) 推出：g 上不变双线性形式空间的维数等于 g 的单分量个数，且所有这些形式都是对称的。
\item
  设 g 为单李代数，M 为向量空间 g 的对偶空间，带有伴随表示的对偶表示，b 为由 \(x \mapsto x_{M}\) 定义的 g 与 M 的半直积。对 g 中的 y、z 与 M 中的 \(y'\)、\(z'\)，写 \(\beta(y + y', z + z') = \langle y, z' \rangle + \langle z, y' \rangle\)。证明 \(\beta\) 是 b 上的非退化不变对称双线性形式，且不与 b 的任何表示相关联。（注意 b 的根基与幂零根基都等于 M。）
\end{enumerate}

\begin{enumerate}
\def\labelenumi{\arabic{enumi}.}
\setcounter{enumi}{18}
\tightlist
\item
  设 \(g\) 为约化李代数，\(U\) 为其包络代数，\(Z\) 为 \(U\) 的中心，\(V\) 为 \(U\) 中由形如 \(uv - vu\)（ \(u \in U, v \in U\) ）的元素生成的子空间。
\end{enumerate}

\begin{enumerate}
\def\labelenumi{(\alph{enumi})}
\item
  U 是 Z 与 V 的直和。（把第 3 节命题 6 应用于表示 \(x \mapsto \mathrm{ad}_U x\)（\(\mathfrak{g}\) 在 U 上的表示），并注意在第 2 节第 7 号定理 1 的推论 4 的分解 \(\mathrm{U} = \sum_{n} \mathrm{U}^n\) 中，诸 \(\mathrm{U}^n\) 在该表示下稳定。）
\item
  设 \(u \mapsto u^{\natural}\) 为 U 到 Z 上平行于 V 的投影。证明 \((uv)^{\natural} = (vu)^{\natural}, (zu)^{\natural} = zu^{\natural}\) 对一切 \(u \in \mathrm{U}, v \in \mathrm{U}, z \in \mathrm{Z}\) 成立。
\item
  设 R 为 U 的双边理想。证明 \(\mathbf{R} = (\mathbf{R} \cap \mathbf{Z}) + (\mathbf{R} \cap \mathbf{V})\)。（为证 R 的元素 r 在 V 中的分量属于 R，归结为 K 代数闭的情形；注意 R 在 U 上的表示 \(\rho\)（它延拓 \(x \mapsto ad_{U}x\)）下稳定；把 r 分解为它在诸 \(U^{n}\) 中的分量；最后，把《代数》第 VIII 章第 4 节第 2 号定理 1 的推论 2 应用于 \(\rho\) 在 \(U^{n}\) 上的限制。）
\item
  设 \(\mathbf{R}'\) 为 \(Z\) 的理想。设 \(\mathbf{R}_1\) 为 \(U\) 中由 \(\mathbf{R}'\) 生成的双边理想。证明 \(\mathbf{R}_1 \cap Z = \mathbf{R}'\)。（利用 (b)。）
\end{enumerate}

\begin{enumerate}
\def\labelenumi{\arabic{enumi}.}
\setcounter{enumi}{19}
\tightlist
\item
  假设 \(\mathbf{K}\) 代数闭。设 \(\mathfrak{g}\) 为李代数，\(\mathfrak{c}\) 为其中心。
\end{enumerate}

\begin{enumerate}
\def\labelenumi{(\alph{enumi})}
\item
  若 \(\mathfrak{g}\) 容许有限维忠实的单表示，则 \(\mathfrak{g}\) 约化且 \(\dim \mathfrak{c} \leqslant 1\)。
\item
  反之，若 g 约化且 dim \(c \leqslant 1\)，则 g 容许有限维忠实的单表示。（若 g 单，用伴随表示。若 \(g = g_{1} \times g_{2}\) 且 \(g_{1}, g_{2}\) 分别容许有限维忠实的单表示 \(\rho_{1}, \rho_{2}\)（在空间 \(M_{1}, M_{2}\) 上），证明
\end{enumerate}

\[
(x _ {1}, x _ {2}) \mapsto \rho_ {1} (x _ {1}) \otimes 1 + 1 \otimes \rho_ {2} (x _ {2})
\]

是半单表示 \(\rho\)（\(g\) 的），再证明 \(g\) -模 \(\mathrm{M}_1 \otimes \mathrm{M}_2\) 的中心化子归结为标量，从而 \(\rho\) 单。当 \(g_1\) 与 \(g_2\) 都半单，或 \(g_1\) 半单且 \(g_2\) 交换时，通过考察 \(\rho\) 的核证明 \(\rho\) 忠实（该核是 \(g_1 \times g_2\) 的理想）。

\begin{enumerate}
\def\labelenumi{\arabic{enumi}.}
\setcounter{enumi}{20}
\item
  设 V 为 K 上有限维向量空间，维数为 n \textgreater{} 1。证明 \(\mathfrak{sl}(V)\) 单。（归结为 K 代数闭的情形。设 \(\mathfrak{sl}(V) = a \times b\)，其中 dim a = a \textgreater{} 0，dim b = b \textgreater{} 0。设 A（相应地 B）为由 1 与 a（相应地 b）生成的结合代数。于是 V 可视为单 \((\mathrm{A} \otimes_{\mathrm{K}} \mathrm{B})\) -模。因而存在 K 上有限维 p 的 A-模 P 与 K 上有限维 q 的 B-模 Q，使得 V \((\mathrm{A} \otimes_{\mathrm{K}} \mathrm{B})\) -同构于 \(P \otimes_{K} Q\)（《代数》第 VIII 章第 7 节第 7 号命题 8 及第 4 号定理 2）；P 与 Q 忠实，因而 \(a \leqslant p^{2}\)、\(b \leqslant q^{2}\)。另一方面，\(a + b = n^{2} - 1\) 且 pq = n，于是 \((p^{2} - 1)(q^{2} - 1) \leqslant 2\)，矛盾。）
\item
  \begin{enumerate}
  \def\labelenumii{(\alph{enumii})}
  \tightlist
  \item
    设 \(\mathfrak{g}\) 为任意特征的代数闭域 \(\mathbf{K}\) 上的幂零李代数。设 \(\mathfrak{z}\) 为其中心，\(\rho\) 为 \(\mathfrak{g}\) 的有限维表示，\(\beta\) 为相应的双线性形式。证明 \(\mathfrak{z}\cap \mathcal{D}\mathfrak{g}\) 与 \(\mathfrak{g}\) 关于 \(\beta\) 正交。（利用 Jordan-Hölder 列归结为 \(\rho\) 单的情形。然后注意，对 \(z\in \mathfrak{g}\cap \mathcal{D}\mathfrak{g},\rho (z)\) 是迹为零的标量矩阵。再利用第 4 节练习 22 (b)。）由此推出，若 \(\beta\) 非退化，则 \(\mathfrak{g}\) 交换。（利用第 4 节练习 3 (a)。）
  \end{enumerate}
\end{enumerate}

\begin{enumerate}
\def\labelenumi{(\alph{enumi})}
\setcounter{enumi}{1}
\item
  假设 \(\mathbf{K}\) 的特征为 2。证明对可解李代数 \(\mathfrak{gl}(2,\mathbf{K}) = \mathfrak{g}\)，与恒等表示相应的双线性形式非退化，但中心 \(\mathfrak{z}\) 包含在 \(\mathcal{D}\mathfrak{g}\) 中且 \(\neq \{0\}\)。
\item
  设 g 为 K 上以 \((a, b, c, d, e, f)\) 为基的 6 维李代数，其乘法表为 \([a, b] = -[b, a] = d\)、\([a, c] = -[c, a] = e\)、\([b, c] = -[c, b] = f\)，其余括积均为 0。设 β 为 g 上的双线性形式，满足 β \((c, d) = \beta(d, c) = 1\)、β \((a, f) = \beta(f, a) = 1\)、β \((b, e) = \beta(e, b) = -1\)，且 β 在所论基的有序偶上的其余取值均为 0。证明 β 不变，g 幂零，\(\mathfrak{z} = \mathcal{D}g \neq \{0\}\)，且 β 非退化。
\end{enumerate}

\begin{enumerate}
\def\labelenumi{\arabic{enumi}.}
\setcounter{enumi}{22}
\tightlist
\item
  设 \(\mathfrak{g}\) 为任意特征的域 \(\mathbf{K}\) 上的 3 维非交换李代数。
\end{enumerate}

\begin{enumerate}
\def\labelenumi{(\alph{enumi})}
\item
  若 \(\mathfrak{g}\) 的中心 \(\mathfrak{z} \neq \{0\}\)，则 dim \(\mathfrak{z} = 1\)（第 1 节练习 2）且 dim \(\mathcal{D}_{\mathfrak{g}} = 1\)。若 \(\mathfrak{z} \neq \mathcal{D}_{\mathfrak{g}}, \mathfrak{g}\) 是 \(\mathfrak{z}\) 与 2 维非交换代数的乘积。若 \(\mathfrak{z} = \mathcal{D}_{\mathfrak{g}}, \mathfrak{g}\) 是 3 维非交换幂零代数（第 4 节练习 9 (a)）。
\item
  若 \(\mathfrak{z} = \{0\}\) 但 \(\mathfrak{g}\) 中存在 2 维交换子代数，则该子代数唯一且等于 \(\mathcal{D}_{\mathfrak{g}}\)；对一切 \(x \notin \mathcal{D}_{\mathfrak{g}}\)，\(u\)（\(ad x\) 在 \(\mathcal{D}_{\mathfrak{g}}\) 上的限制）是该向量空间的一个双射，且在相差一个乘法常数下唯一确定。反之，每个自同构 \(u\)（2 维向量空间 \(\mathrm{Ka} + \mathrm{Kb}\) 的）都在 \(\mathfrak{g} = \mathrm{Ka} + \mathrm{Kb} + \mathrm{Kc}\) 上由条件 \([a, b] = 0\)、\([c, a] = u(a)\)、\([c, b] = u(b)\) 确定一个李代数结构。两个如此由自同构 \(u_1, u_2\) 定义的李代数同构的充分必要条件是 \(u_1\) 与 \(u_2\) 的矩阵在相差一个标量因子下相似。
\item
  假设 g 中不存在维数 \textgreater1 的交换子代数。证明此时存在 \(a \in g\)，使 \(a \notin (\operatorname{ad} a)g\)（假设某元素 x 不具有此性质，利用 Jacobi 恒等式证明另一元素具有所需性质）。于是存在 g 的基 \((a, b, c)\)，使 \([a, b] = c\)、\([a, c] = \beta b\)、\([b, c] = \gamma a\)，其中 \(\beta \neq 0\) 且 \(\gamma \neq 0\)，且 g 单。设 \(\phi\) 为 g 到外幂 \(\bigwedge^{2}\mathfrak{g}^{*}\)（\(\mathfrak{g}\) 的对偶空间的外幂）上的典范同构，它在相差一个可逆常数因子下确定（《代数》第 III 章第 8 节第 5 号）；设 \(u\) 为 \(\mathfrak{g}\) 到 \(\mathfrak{g}^{*}\) 的、由条件 \(\langle [x,y],u(z)\rangle = \langle x\wedge y,\phi (z)\rangle\) 定义的线性映射；双线性形式 \(\Phi (x,y) = \langle x,u(y)\rangle\)（在 \(\mathfrak{g}\times \mathfrak{g}\) 上）对称且非退化，且 \(2\Phi\) 是 \(\mathfrak{g}\) 的 Killing 形式，相差一个 \(\neq 0\) 的因子。K 上两个 3 维单李代数同构的充分必要条件是相应的双线性形式在相差一个 \(\neq 0\) 的常数因子下等价。
\item
  若 \(\mathbf{K}\) 的特征不是 2，则 (c) 中定义的单代数 \(\mathfrak{g}\) 同构于李代数 \(\mathfrak{o}(\Phi)\)（即形式正交群 \(\mathbf{G}(\Phi)\) 的李代数）（第 1 节练习 26 (a)）。要在 \(\mathfrak{g}\) 中存在向量 \(x\)，使 ad \(x\) 容许与 \(x\) 不成比例的特征向量，其充分必要条件是 \(\Phi\) 的指标 \(>0\)；此时 \(\mathfrak{g}\) 容许一个基 \((a, b, c)\)，使 \([a, b] = b\)、\([a, c] = -c\)、\([b, c] = a\)。
\item
  若 K 的特征为 2，证明 g 上不存在 2-映射，因而 g 不是某个形式群的李代数。证明 g 容许非内导子的导子。
\end{enumerate}

¶ 24. 设 K 为任意特征 p 的域。~我们再次采用定义 1。

\begin{enumerate}
\def\labelenumi{(\alph{enumi})}
\tightlist
\item
  证明除非 \(p = n = 2\)，\(\mathfrak{gl}(n, \mathbf{K})\) 的仅有的非平凡理想
\end{enumerate}

是 1 维中心 \(\mathfrak{z}\)（以 \(\sum_{i=1}^{n} E_{ii}\) 为基）以及由迹为 0 的矩阵组成的李代数 \(\mathfrak{sl}(n, \mathrm{K})\)。（除所指出的例外情形外，注意若理想 \(\mathfrak{a}\) 包含某个 \(E_{ij}\)，则它必然包含 \(\mathfrak{sl}(n, \mathrm{K})\)。若 \(\mathfrak{a}\) 包含一个不属于 \(\mathfrak{b}\) 的元素，则用至多四个适当选取的元素 \(E_{ij}\) 去乘该元素，即得到某个 \(E_{ij}\) 的非零倍数。）若 \(n\) 不是 \(p\) 的倍数，证明 \(\mathfrak{gl}(n, \mathrm{K})\) 是 \(\mathfrak{z}\) 与 \(\mathfrak{sl}(n, \mathrm{K})\) 的直和，且 \(\mathfrak{sl}(n, \mathrm{K})\) 单。反之，若 \(n\) 是 \(p\) 的倍数且 \(n > 2\)，证明 \(\mathfrak{sl}(n, \mathrm{K}) / \mathfrak{z}\) 单（同样的方法）。

\begin{enumerate}
\def\labelenumi{(\alph{enumi})}
\setcounter{enumi}{1}
\tightlist
\item
  证明当 \(n\) 是 \(p\) 的倍数且 \(n > 2\) 时，\(\mathfrak{gl}(n, \mathbf{K}) / \mathfrak{sl}(n, \mathbf{K})\) 的根基为 \(\{0\}\)，但它有交换的商 \(\mathfrak{gl}(n, \mathbf{K}) / \mathfrak{sl}(n, \mathbf{K})\)。
\end{enumerate}

¶ 25. 设 K 为任意特征 p 的域。~设 \(\mathfrak{sp}(2n, \mathrm{K})\) 为 K 上 2n 个变量的形式辛群的李代数（第 1 节练习 26 (c)）。证明该李代数（它与 \(\mathfrak{gl}(2n, \mathrm{K})\) 的一个子代数恒同）具有由下列元素组成的基

\[
\begin{array}{c} H _ {i} = E _ {i i} - E _ {i + n, i + n} \quad (1 \leqslant i \leqslant n), \\ F _ {i j} = E _ {i j} - E _ {j + n, i + n} \quad (1 \leqslant i, j \leqslant n, i \neq j), \\ G _ {i j} ^ {\prime} = E _ {i, j + n} + E _ {j, i + n}, \qquad G _ {i j} ^ {\prime \prime} = E _ {i + n, j} + E _ {j + n, i} \quad (1 \leqslant i <   j \leqslant n), \\ E _ {i, i + n} \quad \text { and } \quad E _ {i + n, i} \quad (1 \leqslant i \leqslant n). \end{array}
\]

\begin{enumerate}
\def\labelenumi{(\alph{enumi})}
\item
  证明若 \(p \neq 2\) 且 \(n \geqslant 1\)，则代数 \(\mathfrak{sp}(2n, \mathrm{K})\) 单（当 \(n = 1\) 时，\(\mathfrak{sp}(2, \mathrm{K}) = \mathfrak{sl}(2, \mathrm{K})\)）。（与练习 24 同法。）
\item
  若 \(p = 2\) 且 \(n \geqslant 3\)，证明 \(H_{i}, F_{ij}, G_{ij}'\) 与 \(G_{ij}''\) 构成理想 \(\mathfrak{a}\)，维数为 \(n(2n - 1)\)；\(\mathfrak{a}\) 中形如
\end{enumerate}

\[
\sum_ {i = 1} ^ {n} \lambda_ {i} H _ {i} + \sum_ {i \neq j} \alpha_ {i j} F _ {i j} + \sum_ {i <   j} \beta_ {i j} G _ {i j} ^ {\prime} + \sum_ {i <   j} \gamma_ {i j} G _ {i j} ^ {\prime \prime}
\]

且满足 \(\sum_{i=1}^{n} \lambda_i = 0\) 的元素构成理想 \(b \subset a\)，维数为 \(n(2n - 1) - 1\)，而 \(\sum_{i=1}^{n} H_i\) 的倍数构成一个理想 \(c\)，即 \(\mathfrak{sp}(2n, K)\) 的中心；\(b, c\) 是 \(\mathfrak{a},\mathfrak{b} = \mathcal{D}(\mathfrak{sp}(2n,\mathbf{K}))\) 仅有的理想，且 \(\mathfrak{b} / \mathfrak{c}\) 单。（同法。）\(\mathfrak{sp}(4,\mathbf{K})\) 在 \(\mathbf{K}\) 的特征为 2 时的理想是什么？

¶ 26. 设 K 为特征 ≠2 的域，Φ 为 K 上 n 维向量空间上的非退化对称双线性形式。

\begin{enumerate}
\def\labelenumi{(\alph{enumi})}
\tightlist
\item
  证明若 \(n \geqslant 5\)，则李代数 \(\mathfrak{o}(\Phi)\) 单。（通过扩张基域，可归结为 \(\Phi\) 的指标为 \(m = [n/2]\) 的情形。例如当 \(n = 2m\) 为偶数时，\(\mathfrak{o}(\Phi)\) 具有由下列元素组成的基
\end{enumerate}

\[
\begin{array}{c} H _ {i} = E _ {i i} - E _ {i + m, i + m}, \quad F _ {i j} = E _ {i j} - E _ {j + m, i + m} \quad (1 \leqslant i, j \leqslant m, i \neq j), \\ G _ {i j} ^ {\prime} = E _ {i, j + m} - E _ {j, i + m} \quad \text { and } \quad G _ {i j} ^ {\prime \prime} = E _ {i + m, j} - E _ {j + m, i} \quad (1 \leqslant i <   j \leqslant m); \end{array}
\]

然后仿照练习 24 与 25 论证。当 \(n = 2m + 1\) 为奇数时同样处理。）

\begin{enumerate}
\def\labelenumi{(\alph{enumi})}
\setcounter{enumi}{1}
\tightlist
\item
  设 \(\Delta\) 为 \(\Phi\) 关于某个基的判别式。设 \(n = 4\)。证明若 \(\Delta\) 在 \(\mathbf{K}\) 中不是平方元，则 \(\mathfrak{o}(\Phi)\) 单。若 \(\Delta\) 是平方元，则 \(\mathfrak{o}(\Phi)\) 是两个互相同构的 3 维单李代数的乘积。（利用《代数》第 IX 章第 9 节练习 16 中描述的 \(\mathbf{G}(\Phi)\) 的结构。）由此给出一个单李代数在扩张基域后变为非单的例子。
\end{enumerate}

\begin{enumerate}
\def\labelenumi{\arabic{enumi}.}
\setcounter{enumi}{26}
\tightlist
\item
  \begin{enumerate}
  \def\labelenumii{(\alph{enumii})}
  \tightlist
  \item
    证明在有限维李 \(p\) -代数中，根基是 \(p\) -理想。（利用第 1 节练习 22 (c)。）
  \end{enumerate}
\end{enumerate}

\begin{enumerate}
\def\labelenumi{(\alph{enumi})}
\setcounter{enumi}{1}
\tightlist
\item
  李 \(p\) -代数 \(\mathfrak{g}\) 的根基为零的充分必要条件是 \(\mathfrak{g}\) 不含交换 \(p\) -理想 \(\neq \{0\}\)。（利用第 1 节练习 22 (c)。）
\end{enumerate}

\exercisesfor{7}

第 7 节的约定除非另有说明均有效。

\begin{enumerate}
\def\labelenumi{\arabic{enumi}.}
\tightlist
\item
  \begin{enumerate}
  \def\labelenumii{(\alph{enumii})}
  \tightlist
  \item
    李 K-代数幂零的充分必要条件是它同构于某个代数 \(n(n, \mathbf{K})\) 的一个子代数。
  \end{enumerate}
\end{enumerate}

\begin{enumerate}
\def\labelenumi{(\alph{enumi})}
\setcounter{enumi}{1}
\tightlist
\item
  假设 K 代数闭。李 K-代数可解的充分必要条件是它同构于某个代数 \(t(n, \mathbf{K})\) 的一个子代数。
\end{enumerate}

\begin{enumerate}
\def\labelenumi{\arabic{enumi}.}
\setcounter{enumi}{1}
\item
  设 g 为李代数，n 为其最大幂零理想。则存在有限维向量空间 V 以及 g 到 \(\mathfrak{sl}(V)\) 的一个子代数上的同构，它把 n 的每个元素映为 V 的幂零自同态。（利用 Ado 定理与第 1 节练习 5。）
\item
  设 g 为李代数，U 为其包络代数。
\end{enumerate}

\begin{enumerate}
\def\labelenumi{(\alph{enumi})}
\item
  证明对一切 \(u \in \mathrm{U}\)（ \(u \neq 0\) ），存在一个有限维表示 \(\pi\)（\(\mathrm{U}\) 上的），使 \(\pi(u) \neq 0\)。（利用练习 2 与第 3 节练习 5 (b)。）
\item
  由 (a) 推出：若 g 半单，则存在无穷多个互不等价的 g 的有限维单表示。（设 \(\rho_{1},\ldots,\rho_{n}\) 为 g 的有限维单表示，\(N_{1},\ldots,N_{n}\) 为相应 U-模的零化理想，\(N=\bigcap_{i=1}^{n}N_{i}\)。证明 N 在 U 中有余限维。设 \(n\in N,n\neq0\)。对 n 应用 (a)。）
\end{enumerate}

\begin{enumerate}
\def\labelenumi{\arabic{enumi}.}
\setcounter{enumi}{3}
\tightlist
\item
  设 g 为李代数，a 为 g 的次不变子代数，\(\mathfrak{r}(a)\) 为 a 的根基，\(\rho\) 为 a 的有限维表示。\(\rho\) 容许到 g 的有限维扩张的充分必要条件是 \(\mathcal{D}_{\mathfrak{g}} \cap \mathfrak{r}(a)\) 包含在 \(\rho\) 的最大幂零性理想中。（为证充分性，如下进行：设 \(g = g_{0} \supset g_{1} \supset \cdots \supset g_{k} = a\) 为 g 的具有第 6 节练习 11 (b) 中诸性质的合成列。用归纳法证明存在扩张 \(\rho_{i}\)（\(\rho\) 到 \(g_{0}\) 的扩张），其最大幂零性理想包含 \(\mathcal{D}_{\mathfrak{g}} \cap \mathfrak{r}(g_{i})\)。在第 6 节练习 11 (b) 的记号下，由 \(\rho_{i+1}\) 到 \(\rho_{i}\) 的过渡由定理 1 得出：当 \(b_{i}\) 单，或当 \(b_{i}\) 为 1 维且
\end{enumerate}

\[
\mathscr {D} \mathfrak {g} \cap \mathfrak {r} (\mathfrak {g} _ {i + 1}) = \mathscr {D} \mathfrak {g} \cap \mathfrak {r} (\mathfrak {g} _ {i}).
\]

时成立。当 \(\mathcal{D}\mathfrak{g} \cap \mathfrak{r}(\mathfrak{g}_{i+1}) \neq \mathcal{D}\mathfrak{g} \cap \mathfrak{r}(\mathfrak{g}_i)\) 时，可设 \(\mathfrak{h}_i \subset \mathcal{D}\mathfrak{g} \cap \mathfrak{r}(\mathfrak{g}_i)\)。于是 \(\mathfrak{r}(\mathfrak{g}_i) = \mathfrak{r}(\mathfrak{g}) \cap \mathfrak{g}_i\)（第 5 节练习 10），从而可应用定理 1。）

¶ 5. 设 K 为特征 p \textgreater{} 0 的域，g 为 K 上有限维 n 李代数，U 为 g 的包络代数。

\begin{enumerate}
\def\labelenumi{(\alph{enumi})}
\item
  \(p\) -多项式（\(\mathbf{K}\) 上、一个不定元的）是指 \(\mathrm{K}[\mathrm{X}]\) 中仅有的 \(\neq 0\) 项都是 \(p\) 的方幂次数的多项式。 证明对每个多项式 \(f(\mathbf{X}) \neq 0\)（在 \(\mathrm{K}[\mathrm{X}]\) 中），存在 \(g(\mathbf{X}) \in \mathrm{K}[\mathbf{X}]\) ，使 \(f(\mathbf{X})g(\mathbf{X})\) 是 \(p\) -多项式（考察诸 \(\mathbf{X}^{p_k}\) 被 \(f(\mathbf{X})\) 作欧几里得除法所得的余式）。
\item
  证明对每个元素 \(z \in \mathfrak{g}\)，存在非零多项式 \(f(\mathbf{X}) \in \mathbf{K}[\mathbf{X}]\)，使 \(f(z)\) 属于 U 的中心 C。（考察自同态 ad \(z\) 的最小多项式，对此多项式应用 (a)，并利用第 1 节练习 19 的公式 (1)。）
\item
  设 \((e_i)_{1\leqslant i\leqslant n}\) 为 \(\mathfrak{g}\) 的基。对每个 \(i\)，设 \(f_{i}\neq 0\) 为 \(\mathrm{K}[\mathrm{X}]\) 中使 \(f_{i}(e_{i})\in \mathrm{C}\) 的多项式；设 \(d_{i}\) 为其次数，\(\mathrm{L}\) 为 \(\mathrm{U}\) 中由诸 \(y_{i} = f_{i}(e_{i})\) 生成的双边理想。证明模 \(\mathrm{L}\) 的、元素 \(e_1^{\alpha_1}\ldots e_n^{\alpha_n}\)（其中 \(0\leqslant \alpha_{i} < d_{i}\)）的类构成 \(\mathrm{U / L}\) 的基。
\item
  证明 \(\mathfrak{g}\) 容许有限维忠实的线性表示（选取诸 \(f_{i}\)，使 \(d_{i} > 1\) 对一切 \(i\) 成立）。
\item
  证明若 \(g \neq \{0\}\)，则存在 g 的非半单有限维线性表示。（设诸 \(f_{t}\) 的次数为 \(d_{t} > 1\)，在 L 的定义中把诸 \(f_{t}\) 换成 \(f_{t}^{2}\)，并注意此时 U/L 的中心中存在 \(\neq 0\) 的幂零元素，因而 U/L 不是半单的。）
\end{enumerate}


\chapter{自由李代数}

本章中，†字母 K 表示一个非零交换环。K 的单位元素记作 1。除非另有说明，所有余代数、代数与双代数、所有模以及所有张量积都是在 K 上的。

从第 6 节起，将假设 K 是特征 0 的域。

\section{李代数的包络双代数}

本节中，\(\mathfrak{g}\) 表示 \(\mathbf{K}\) 上的李代数，\(\mathrm{U}(\mathfrak{g})\) 或简称 \(\mathrm{U}\) 表示它的包络代数（第 I 章，第 2 节，第 1 号），\(\sigma\) 表示 \(\mathfrak{g}\) 到 \(\mathrm{U}(\mathfrak{g})\) 的典范映射（同处），\((\mathrm{U}_n)_{n\geqslant 0}\) 表示 \(\mathrm{U}\) 的典范滤过（同处，第 6 号）。

\subsection{余代数的本原元素}

本号中我们考察一个余代数 E（《代数》，第 III 章，第 11 节，第 1 号），它带有余积

\[
c \colon \mathrm{E} \rightarrow \mathrm{E} \otimes \mathrm{E}
\]

和余单位 \(\varepsilon\)（同处，第 2 号）。回忆 \(\varepsilon\) 是 K-模 E 上的线性形式，满足（在 \(\mathrm{E} \otimes \mathrm{K}\) 与 \(\mathrm{K} \otimes \mathrm{E}\) 和 E 的典范等同下）：

\[
\mathrm{Id} _ {\mathrm{E}} = (\varepsilon \otimes \mathrm{Id} _ {\mathrm{E}}) \circ c = (\mathrm{Id} _ {\mathrm{E}} \otimes \varepsilon) \circ c.
\]

设 \(E^{+}\) 表示 \(\varepsilon\) 的核，并设 u 为 E 的元素，使得

\[
c (u) = u \otimes u \quad \text { and } \quad \varepsilon (u) = 1.
\]

K-模 E 是 \(E^{+}\) 与子模 K.u 的直和，后者是以 u 为基的自由模；设 \(\pi_{u}: E \to E^{+}\) 与 \(\eta_{u}: E \to K.u\) 为与这一分解相伴的射影。那么

\[
\pi_ {u} (x) = x - \varepsilon (x). u, \quad \eta_ {u} (x) = \varepsilon (x). u.\tag{1}
\]

定义 1. 元素 \(x\)（\(\mathbf{E}\) 的元素）称为 \(u\) -本原的，如果

\[
c (x) = x \otimes u + u \otimes x.\tag{2}
\]

\(u\) -本原的诸元素（\(\mathbf{E}\) 的）构成 \(\mathbf{E}\) 的一个子模，记作 \(\mathrm{P}_{u}(\mathrm{E})\)。

命题 1. 每个 \(u\) -本原元素（\(\mathbf{E}\) 的）都属于 \(\mathbf{E}^{+}\)。

\begin{enumerate}
\def\labelenumi{(\arabic{enumi})}
\setcounter{enumi}{1}
\tightlist
\item
  蕴含 \(x = \varepsilon(x).u + \varepsilon(u).x = \varepsilon(x).u + x\)，由此 \(\varepsilon(x) = 0\)。
\end{enumerate}

注. 若 \(x \in \mathrm{E}\) 且 \(c(x) = x' \otimes u + u \otimes x''\)，其中 \(x', x''\) 属于 \(\mathrm{E}^+\)，则 \(x = \varepsilon(x')\) . \(u + \varepsilon(u)\) . \(x'' = x''\)；同理 \(x = x'\)，于是 \(x\) 是 \(u\) -本原的。

对一切 \(x \in \mathrm{E}^{+}\)，我们记

\[
c _ {u} ^ {+} (x) = c (x) - x \otimes u - u \otimes x.\tag{3}
\]

命题 2. 我们有

\[
\left(\pi_ {u} \otimes \pi_ {u}\right) \circ c = c _ {u} ^ {+} \circ \pi_ {u}.\tag{4}
\]

设 \(x\) 属于 E；那么

\[
\begin{array}{r l} (\pi_ {u} \otimes \pi_ {u}) (c (x)) & = ((1 - \eta_ {u}) \otimes (1 - \eta_ {u})) (c (x)) \\ & = c (x) - (1 \otimes \eta_ {u}) (c (x)) - (\eta_ {u} \otimes 1) (c (x)) + (\eta_ {u} \otimes \eta_ {u}) (c (x)). \end{array}
\]

由于 \(\varepsilon\) 是 E 的余单位，

\[
(1 \otimes \eta_ {u}) (c (x)) = x \otimes u, \quad (\eta_ {u} \otimes 1) (c (x)) = u \otimes x
\]

由此

\[
\left(\eta_ {u} \otimes \eta_ {u}\right) (c (x)) = \left(\eta_ {u} \otimes 1\right) \left(\left(1 \otimes \eta_ {u}\right) (c (x))\right) = \varepsilon (x). u \otimes u;
\]

由此我们得出

\[
\left(\pi_ {u} \otimes \pi_ {u}\right) (c (x)) = c (x) - x \otimes u - u \otimes x + \varepsilon (x). u \otimes u.
\]

另一方面，

\[
c _ {u} ^ {+} \left(\pi_ {u} (x)\right) = c (x) - x \otimes u - u \otimes x + \varepsilon (x). u \otimes u,
\]

由此得公式 (4)。

由于 \(\mathrm{E}^{+}\) 是 \(\mathrm{E}\) 的直因子子模，\(\mathrm{E}^{+} \otimes \mathrm{E}^{+}\) 可等同于 \(\mathrm{E} \otimes \mathrm{E}\) 的一个直因子子模。在这一等同下，\(\pi_{u} \otimes \pi_{u}\) 是 \(\mathrm{E} \otimes \mathrm{E}\) 到 \(\mathrm{E}^{+} \otimes \mathrm{E}^{+}\) 上的射影。由公式 (4)，\(c_{u}^{+}\) 把 \(\mathrm{E}^{+}\) 映入 \(\mathrm{E}^{+} \otimes \mathrm{E}^{+}\)，且 \(\pi_{u}\) 是余代数 \((\mathrm{E}, c)\) 到余代数 \((\mathrm{E}^{+}, c_{u}^{+})\) 的态射。

命题 3. 若余代数 \((\mathrm{E}, c)\) 是余结合的（相应地，余交换的）（《代数》，第 III 章，第 11 节，第 2 号），则余代数 \((\mathrm{E}^{+}, c_{\mathbf{u}}^{+})\) 亦然。这由下述引理得出：

引理 1. 设 \(\pi: \mathrm{E} \to \mathrm{E}'\) 为满射余代数态射。若 \(\mathrm{E}\) 是余结合的（相应地，余交换的），则 \(\mathrm{E}'\) 亦然。

设 B 为结合 K-代数；映射 \(f \mapsto f \circ \pi\) 是 \(\mathrm{Hom}_{\mathbb{K}}(\mathrm{E}', \mathrm{B})\) 到 \(\mathrm{Hom}_{\mathbb{K}}(\mathrm{E}, \mathrm{B})\) 的单射代数同态。于是只需应用《代数》第 III 章第 11 节第 2 号的命题 1（相应地，命题 2）。

\subsection{双代数的本原元素}

设 E 为双代数（《代数》，第 III 章，第 11 节，第 4 号），c 为其余积，ε 为其余单位，1 为其单位元素。由于 ε(1) = 1 且 c(1) = 1 ⊗ 1，前一号的结果可在 u = 1 下应用。E 的 1-本原元素（第 1 号，定义 1）简称为本原元素（参看~《代数》，第 III 章，第 11 节，第 8 号），即 E 中满足下式的元素 x：

\[
c (x) = x \otimes 1 + 1 \otimes x.\tag{5}
\]

我们用 \(\pi\)、\(\eta\)、P(E)、\(c^{+}\) 简记 \(\pi_{1}\)、\(\eta_{1}\)、P \(_{1}\) (E)、\(c_{1}^{+}\)。

命题 4. 本原元素的集合 \(\mathrm{P(E)}\)（\(\mathbf{E}\) 的）是 \(\mathbf{E}\) 的李子代数。

若 \(x, y\) 属于 \(\mathrm{P(E)}\)，则

\[
\begin{array}{r l} c (x y) & = c (x) c (y) = (x \otimes 1 + 1 \otimes x) (y \otimes 1 + 1 \otimes y) \\ & = x y \otimes 1 + 1 \otimes x y + x \otimes y + y \otimes x, \end{array}
\]

由此

\[
c ([ x, y ]) = [ x, y ] \otimes 1 + 1 \otimes [ x, y ].
\]

命题 5. 设 \(f: \mathrm{E} \to \mathrm{E}'\) 为双代数态射。若 \(x\) 是 \(\mathrm{E}\) 的本原元素，则 \(f(x)\) 是 \(\mathrm{E}'\) 的本原元素，且 \(f\) 在 \(\mathrm{P(E)}\) 上的限制是李代数同态 \(\mathrm{P}(f): \mathrm{P(E)} \to \mathrm{P(E')}\)。

设 \(c\)（相应地，\(c'\)）为 E（相应地，E'）的余积。由于 \(f\) 是余代数态射，\(c' \circ f = (f \otimes f) \circ c\)，由此

\[
\begin{array}{r l} c ^ {\prime} (f (x)) & = (f \otimes f) (c (x)) = (f \otimes f) (x \otimes 1 + 1 \otimes x) \\ & = f (x) \otimes 1 + 1 \otimes f (x), \end{array}
\]

其中 \(x\) 是本原的。因而 \(f\) 把 \(\mathrm{P(E)}\) 映入 \(\mathrm{P(E')}\)，且 \(f([x,y]) = [f(x),f(y)]\)，因为 \(f\) 是代数同态。

注. (1) 设 \(p\) 为素数，使得 \(p\) . \(1 = 0\) 在 \(\mathbf{K}\) 中成立。二项式公式以及同余关系 \(\binom{p}{i} \equiv 0 \pmod{p}\)（对 \(1 \leqslant i \leqslant p - 1\)）蕴含

P(E) 在映射 \(x \mapsto x^p\) 下稳定。

\begin{enumerate}
\def\labelenumi{(\arabic{enumi})}
\setcounter{enumi}{1}
\tightlist
\item
  按定义，图示
\end{enumerate}

\[
0 \longrightarrow \mathrm{P(E)} \longrightarrow \mathrm {E^ {+}} \stackrel {c ^ {+}} {\longrightarrow} \mathrm {E^ {+}} \otimes \mathrm {E^ {+}}
\]

是正合列。若 \(\mathbf{K}'\) 是交换环且 \(\rho: \mathbf{K} \to \mathbf{K}'\) 是环同态，则 \(\rho^*(\mathrm{E}) = \mathrm{E} \otimes_{\mathbb{K}} \mathbf{K}'\) 是 \(\mathbf{K}'\) -双代数，且含入 \(\mathrm{P}(\mathrm{E}) \to \mathrm{E}\) 定义了李 \(\mathbf{K}'\) -代数的同态

\[
\alpha \colon \mathrm{P} (\mathrm{E}) \otimes_ {\mathrm{K}} \mathrm{K} ^ {\prime} \rightarrow \mathrm{P} (\mathrm{E} \otimes_ {\mathrm{K}} \mathrm{K} ^ {\prime}).
\]

若 \(\mathbf{K}'\) 在 \(\mathbf{K}\) 上是平坦的（《交换代数》，第 I 章，第 2 节，第 3 号，定义 2），则由同处可知图示

\[
0 \longrightarrow \mathrm{P} (\mathrm{E}) \otimes_ {\mathrm{K}} \mathrm{K} ^ {\prime} \longrightarrow \mathrm{E} ^ {+} \otimes_ {\mathrm{K}} \mathrm{K} ^ {\prime} \xrightarrow {c ^ {+} \otimes_ {\mathrm{K}} \mathrm{Id} _ {\mathrm{K} ^ {\prime}}} \left(\mathrm{E} ^ {+} \otimes_ {\mathrm{K}} \mathrm{K} ^ {\prime}\right) \otimes_ {\mathrm{K} ^ {\prime}} \left(\mathrm{E} ^ {+} \otimes_ {\mathrm{K}} \mathrm{K} ^ {\prime}\right)
\]

是正合列，由此推出 \(\alpha\) 是同构。

\subsection{滤过双代数}

定义 2. 设 E 为带有余积 c 的双代数。~与 E 上双代数结构相容的滤过是指 E 的子模的递增序列 \((\mathrm{E}_{n})_{n\geq0}\)，满足

\[
\begin{array}{c} \mathrm{E} _ {0} = \mathrm{K}. 1, \qquad \mathrm{E} = \bigcup_ {n \geqslant 0} \mathrm{E} _ {n} \\ \mathrm{E} _ {m}. \mathrm{E} _ {n} \subset \mathrm{E} _ {m + n} \quad \text {for} \quad m \geqslant 0, n \geqslant 0 \\ c (\mathrm{E} _ {n}) \subset \sum_ {i + j = n} \mathrm{Im} (\mathrm{E} _ {i} \otimes \mathrm{E} _ {j}) \quad \text {for} n \geqslant 0. \dagger \end{array}\tag{6}
\]

带有与其双代数结构相容的滤过的双代数称为滤过双代数。

例. 设 E 为分次双代数（《代数》，第 III 章，第 11 节，第 4 号，定义 3），\((\mathrm{E}^n)_{n\geqslant 0}\) 为其分次。我们记 \(\mathrm{E}_n = \sum_{i = 0}^{n}\mathrm{E}^i\)。序列 \((\mathrm{E}_n)\) 是与 E 上双代数结构相容的滤过。

命题 6. 设 E 为滤过双代数，\((\mathrm{E}_n)_{n\geqslant 0}\) 为其滤过。对每个整数 \(n\geqslant 0\)，设 \(\mathrm{E}_n^+ = \mathrm{E}_n\cap \mathrm{E}^+\)。那么 \(\mathrm{E_0^+} = \{0\}\) 且

\[
c ^ {+} (\mathrm{E} _ {n} ^ {+}) \subset \sum_ {i = 1} ^ {n - 1} \operatorname{Im} (\mathrm{E} _ {i} ^ {+} \otimes \mathrm{E} _ {n - i} ^ {+}) \quad \text { for } n \geqslant 0. \dagger\tag{7}
\]

由于 \(\mathrm{E}_0 = \mathrm{K}.1, \mathrm{E}_0^+ = 0\)。若 \(x \in \mathrm{E}_n\)，则 \(\pi(x) = x - \varepsilon(x).1\)（公式 (1)），由此 \(\pi(x) \in \mathrm{E}_n^+\)，即 \(\pi(\mathrm{E}_n) \subset \mathrm{E}_n^+\)。于是 \(\pi \otimes \pi\) 把 \(\operatorname{Im}(\mathrm{E}_i \otimes \mathrm{E}_j)\) 映入 \(\operatorname{Im}(\mathrm{E}_i^+ \otimes \mathrm{E}_j^+)\)（对 \(i \geqslant 0, j \geqslant 0\)）。由于 \(c^+ = (\pi \otimes \pi) \circ c\) 在 \(\mathrm{E}^+\) 中成立（第 1 号，命题 2），由 (6) 得

\[
c ^ {+} \left(\mathrm{E} _ {n} ^ {+}\right) \subset \sum_ {i = 0} ^ {n} \operatorname{Im} \left(\mathrm{E} _ {i} ^ {+} \otimes \mathrm{E} _ {n - i} ^ {+}\right) = \sum_ {i = 1} ^ {n - 1} \operatorname{Im} \left(\mathrm{E} _ {i} ^ {+} \otimes \mathrm{E} _ {n - i} ^ {+}\right).
\]

推论. \(\mathbf{E}_1^+\) 的元素都是本原的。

若 \(x \in \mathrm{E}_{1}^{+}\)，则由 (7) 得 \(c^{+}(x) = 0\)，由此得 (5)。

\subsection{李代数的包络双代数}

回忆 g 表示李代数，U 表示其包络代数，连同其典范滤过 \((\mathrm{U}_{n})_{n\geqslant0}\)。

命题 7. 在代数 U 上存在且仅存在一个余积 c，使 U 成为双代数且 \(\sigma(g)\) 的元素都是本原的。双代数 (U, c) 是余交换的；其余单位是线性形式 \(\varepsilon\)，使得 U 的每个元素 \(x\) 的常数项（第 I 章，第 2 节，第 1 号）为 \(\varepsilon(x)\) .1。U 的典范滤过 \((\mathrm{U}_n)_{n\geqslant 0}\) 与这一双代数结构相容。

\begin{enumerate}
\def\labelenumi{(\alph{enumi})}
\tightlist
\item
  设 \(x \in \mathfrak{g}\)；我们记 \(c_0(x) = \sigma(x) \otimes 1 + 1 \otimes \sigma(x) \in \mathrm{U} \otimes \mathrm{U}\)。若 \(x, y\) 属于 \(\mathfrak{g}\)，则
\end{enumerate}

\[
c _ {0} (x) c _ {0} (y) = (\sigma (x) \sigma (y)) \otimes 1 + 1 \otimes (\sigma (x) \sigma (y)) + \sigma (x) \otimes \sigma (y) + \sigma (y) \otimes \sigma (x),
\]

由此

\[
[ c _ {0} (x), c _ {0} (y) ] = c _ {0} ([ x, y ]).
\]

由 U 的泛性质（第 I 章，第 2 节，第 1 号，命题 1），存在且仅存在一个幺代数同态

\[
c \colon \mathrm{U} \to \mathrm{U} \otimes \mathrm{U}
\]

使得 \(c(\sigma(x)) = \sigma(x) \otimes 1 + 1 \otimes \sigma(x)\) 对 \(x \in \mathfrak{g}\) 成立。这就证明了命题 7 中的唯一性断言。

\begin{enumerate}
\def\labelenumi{(\alph{enumi})}
\setcounter{enumi}{1}
\tightlist
\item
  我们证明 \(c\) 是余结合的。线性映射 \(c'\) 与 \(c''\)（\(U\) 到 \(U \otimes U \otimes U\) 的、由下式定义的）
\end{enumerate}

\[
c ^ {\prime} = (c \otimes \mathrm{Id} _ {\mathrm{U}}) \circ c \quad \text { and } \quad c ^ {\prime \prime} = (\mathrm{Id} _ {\mathrm{U}} \otimes c) \circ c
\]

是幺代数同态，且在 \(\sigma (\mathfrak{g})\) 上一致，因为对 \(a\in \sigma (\mathfrak{g})\)

\[
c ^ {\prime} (a) = a \otimes 1 \otimes 1 + 1 \otimes a \otimes 1 + 1 \otimes 1 \otimes a = c ^ {\prime \prime} (a),
\]

由此得结论。

\begin{enumerate}
\def\labelenumi{(\alph{enumi})}
\setcounter{enumi}{2}
\item
  我们证明 \(c\) 是余交换的。设 \(\tau\) 为 \(\mathbf{U} \otimes \mathbf{U}\) 的自同构，使得 \(\tau(a \otimes b) = b \otimes a\) 对 \(a, b\) ∈ \(\mathbf{U}\) 成立。映射 \(\tau \circ c\) 与 \(c\)（都是 \(\mathbf{U}\) 到 \(\mathbf{U} \otimes \mathbf{U}\) 的映射）是幺代数同态，且在 \(\sigma(\mathfrak{g})\) 上一致，由此得结论。
\item
  我们证明 \(\varepsilon\) 是 \(c\) 的余单位。U 到 U 的映射 \((\mathrm{Id}_{\mathbb{U}} \otimes \varepsilon) \circ c\) 与 \((\varepsilon \otimes \mathrm{Id}_{\mathbb{U}}) \circ c\) 是幺代数同态，且与 \(\mathrm{Id}_{\mathbb{U}}\) 在 \(\sigma(\mathfrak{g})\) 上一致。
\item
  我们知道 \(\mathrm{U}_0 = \mathrm{K}.1, \mathrm{U}_n \subset \mathrm{U}_{n+1}, \mathrm{U} = \bigcup_{n \geqslant 0} \mathrm{U}_n\) 以及 \(\mathrm{U}_n. \mathrm{U}_m \subset \mathrm{U}_{n+m}\)（第 I 章，第 2 节，第 6 号）。设 \(a_1, \ldots, a_n\) 属于 \(\sigma(\mathfrak{g})\)。那么
\end{enumerate}

\[
\begin{array}{l} c (a _ {1} \ldots a _ {n}) = \prod_ {i = 1} ^ {n} c (a _ {i}) = \prod_ {i = 1} ^ {n} (a _ {i} \otimes 1 + 1 \otimes a _ {i}) \\ = \sum_ {i = 0} ^ {n} \sum_ {\alpha \in I (i)} (a _ {\alpha (1)} \ldots a _ {\alpha (i)}) \otimes (a _ {\alpha (i + 1)} \ldots a _ {\alpha (n)}), \end{array}\tag{8}
\]

其中 \(I(i)\) 表示 \([1, n]\) 的、在每个区间 \([1, i]\) 与 \([i + 1, n]\) 上都递增的置换的集合。由于 \(U_{n}\) 是由 \(\sigma(g)\) 中至多 n 个元素的乘积生成的 K-模，公式 (8) 蕴含滤过 \((\mathrm{U}_{n})\) 与 \((\mathrm{U}, c)\) 的双代数结构相容。

定义 3. 双代数 \((\mathrm{U}, c)\) 称为李代数 \(\mathfrak{g}\) 的包络双代数。

命题 8. 设 E 为双代数，其余积记作 \(c_{\mathbb{E}}\)，并设 h 为 g 到 P(E) 的李代数同态（第 2 号，命题 4）。使对一切 x ∈ g 有 f(σ(x)) = h(x) 的幺代数同态 f: U → E 是双代数态射。

\[
(f \otimes f) \circ c = c _ {\mathrm{E}} \circ f.
\]

\[
a \in \sigma (\mathfrak {g})
\]

\[
(f \otimes f) (c (a)) = f (a) \otimes 1 + 1 \otimes f (a) = c _ {\mathbb {E}} (f (a))
\]

因为 \(f(a) \in \mathrm{P}(\mathrm{E})\)。类似地，若 \(\varepsilon_{\mathbb{E}}\) 是 \(\mathrm{E}\) 的余单位，则 \(\varepsilon_{\mathbb{E}} \circ f\) 是 \(\mathrm{U} \to \mathrm{K}\) 的幺代数同态，它在 \(\sigma(\mathfrak{g})\) 上为零（第 1 号，命题 1），因而与 \(\varepsilon\) 一致。

由命题 5 与命题 8 可知，映射 \(f \mapsto f \circ \sigma\) 给出双代数同态 \(U(\mathfrak{g}) \to E\) 与李代数同态 \(\mathfrak{g} \to P(E)\) 之间的一一对应。

推论. 设 \(\mathfrak{g}_i\)（\(i = 1,2\)）为李代数，\(\mathrm{U}(\mathfrak{g}_i)\) 为其包络双代数，\(\sigma_i\colon \mathfrak{g}_i\to \mathrm{U}(\mathfrak{g}_i)\) 为典范映射。对每个李代数同态 \(h\colon \mathfrak{g}_1\to \mathfrak{g}_2\)，幺代数同态 \(\mathrm{U}(h)\colon \mathrm{U}(\mathfrak{g}_1)\to \mathrm{U}(\mathfrak{g}_2)\)（它使 \(\mathrm{U}(h)\circ \sigma_1 = \sigma_2\circ h\)，第 I 章，第 2 节，第 1 号）是双代数态射。

\subsection{特征 0 时余代数 U(g) 的结构}

本号中，将假设 K 是特征 0 的域。

设 \(S(\mathfrak{g})\) 为向量空间 \(\mathfrak{g}\) 的对称代数，\(c_{\mathrm{s}}\) 为其余积（《代数》，第 III 章，第 11 节，第 1 号，例 6），\(\eta\) 为向量空间 \(\mathsf{S}(\mathfrak{g})\) 到向量空间 U 上的典范同构（第 I 章，第 2 节，第 7 号）。回忆若 \(x_{1},\ldots ,x_{n}\) 属于 \(\mathfrak{g}\)，则

\[
\eta (x _ {1} \dots x _ {n}) = \frac {1}{n !} \sum_ {\tau \in \mathfrak {S} _ {n}} \sigma (x _ {\tau (1)}) \dots \sigma (x _ {\tau (n)}).\tag{9}
\]

特别地，对 \(x \in g\) 与 \(n \geqslant 0\)，

\[
\eta (x ^ {n}) = \sigma (x) ^ {n}.\tag{10}
\]

注意，由《代数》第 III 章第 6 节第 1 号注 3，\(\eta\) 是 \(\mathbb{S}(\mathfrak{g})\) 到 U 的满足条件 (10) 的唯一线性映射。

命题 9. 对每个整数 \(n \geqslant 0\)，设 \(\mathbf{U}^n\) 为 \(\mathbf{U}\) 的由诸 \(\sigma(x)^n\)（\(x \in \mathfrak{g}\)）生成的向量子空间。

\begin{enumerate}
\def\labelenumi{(\alph{enumi})}
\tightlist
\item
  序列 \((\mathbf{U}^n)_{n\geqslant 0}\) 是向量空间 \(\mathbf{U}\) 的与其余代数结构相容的分次。
\end{enumerate}

给 U 赋以分次 (U \(^{n}\))。

\begin{enumerate}
\def\labelenumi{(\alph{enumi})}
\setcounter{enumi}{1}
\tightlist
\item
  典范映射 \(\eta: \mathbf{S}(\mathfrak{g}) \to \mathbf{U}\) 是分次余代数的同构。设 \(x \in \mathfrak{g}\)，\(n \in \mathbf{N}\)。那么
\end{enumerate}

\[
c _ {\mathrm{S}} (x ^ {n}) = c _ {\mathrm{S}} (x) ^ {n} = (x \otimes 1 + 1 \otimes x) ^ {n} = \sum_ {i = 0} ^ {n} {\binom {n} {i}} x ^ {i} \otimes x ^ {n - i}\tag{11}
\]

因为 \(c_{\mathbb{S}}\) 是代数同态。类似地，由 (10) 得

\[
\begin{array}{l} c (\eta (x ^ {n})) = c (\sigma (x) ^ {n}) = c (\sigma (x)) ^ {n} = (\sigma (x) \otimes 1 + 1 \otimes \sigma (x)) ^ {n} \\ = \sum_ {i = 0} ^ {n} \binom {n} {i} \sigma (x) ^ {i} \otimes \sigma (x) ^ {n - i} = \sum_ {i = 0} ^ {n} \binom {n} {i} \eta (x ^ {i}) \otimes \eta (x ^ {n - i}), \end{array}\tag{12}
\]

由此

\[
(\eta \otimes \eta) (c _ {\mathrm{S}} (x ^ {n})) = c (\eta (x ^ {n})).
\]

由于诸 \(x^n\)（对 \(x \in \mathfrak{g}\) 与 \(n \in \mathbf{N}\)）生成向量空间 \(\mathsf{S}(\mathfrak{g})\)，故 \((\eta \otimes \eta) \circ c_{\mathbf{S}} = c \circ \eta\)，从而 \(\eta\) 是余代数同构。

另一方面，公式 (10) 表明 \(\eta(\mathbb{S}^n(\mathfrak{g})) = \mathrm{U}^n\)，注意到 \(\mathbb{S}(\mathfrak{g})\) 的分次与其余代数结构相容这一事实，这就完成了 (a) 与 (b) 的证明。

U 的分次 \((\mathrm{U}^{n})_{n\geqslant0}\) 称为典范分次。

推论. 典范映射 \(\sigma\) 给出 \(g\) 到李代数 \(\mathrm{P}(\mathbf{U})\)（\(\mathbf{U}\) 的本原元素组成的李代数）上的同构。

由于 \(c^{+}\) 是次数 0 的分次同态，

\[
\mathrm{P} (\mathrm{U}) = \sum_ {n \geqslant 1} (\mathrm{P} (\mathrm{U}) \cap \mathrm{U} ^ {n}).
\]

只需证明：若 \(n > 1\) 且 \(a \in \mathrm{U}^n\) 是本原的，则 \(a = 0\)。今 \(a\) 可写成 \(\sum_{i} \lambda_i a_i^n\)，其中 \(\lambda_i \in \mathbf{K}\)，\(a_i \in \sigma(g)\)。由 (12)，双次数 \((1, n - 1)\) 的项在 \(c^+(a)\) 中为 \(n \sum_{i} \lambda_i a_i \otimes a_i^{n-1}\)。于是 \(\sum_{i} \lambda_i a_i \otimes a_i^{n-1} = 0\)。若 \(\mu: \mathrm{U} \otimes \mathrm{U} \to \mathrm{U}\) 是由 \(\mathrm{U}\) 上乘法所定义的线性映射，则

\[
a = \sum_ {i} \lambda_ {i} a _ {i} ^ {n} = \mu \left(\sum_ {i} \lambda_ {i} a _ {i} \otimes a _ {i} ^ {n - 1}\right) = 0.
\]

注. (1) \(\mathrm{U}_n = \sum_{i=0}^{n} \mathrm{U}^i\)（第 I 章，第 2 节，第 7 号，定理 1 的推论 4）。

\begin{enumerate}
\def\labelenumi{(\arabic{enumi})}
\setcounter{enumi}{1}
\tightlist
\item
  映射 \(\eta\) 是 \(\mathsf{S}(\mathfrak{g})\) 到 U 的分次余代数态射中使得 \(\eta(1) = 1\) 且 \(\eta(x) = \sigma(x)\)（对 \(x \in \mathfrak{g}\)）的唯一态射。因为若 \(\eta'\) 是满足这些条件的态射，我们对 \(n\) 用归纳法证明：\(\eta'(x^n) = \eta(x^n)\)（对 \(x \in \mathfrak{g}\) 与 \(n > 1\)）。由于由 (3) 与 (11) 有 \(c_{\mathbb{S}}^{+}(x^{n}) = \sum_{i=1}^{n-1} \binom{n}{i} x^{i} \otimes x^{n-i}\)，
\end{enumerate}

\[
(\eta \otimes \eta) (c _ {\mathbf {S}} ^ {+} (x ^ {n})) = (\eta^ {\prime} \otimes \eta^ {\prime}) (c _ {\mathbf {S}} ^ {+} (x ^ {n}))
\]

由归纳假设。由此得出 \(c^{+}(\eta(x^{n})) = c^{+}(\eta'(x^{n}))\)；于是 \(\eta(x^{n}) - \eta'(x^{n})\) 是次数 \(n\) 的本原元素，因而为零（命题 9 的推论）。

\begin{enumerate}
\def\labelenumi{(\arabic{enumi})}
\setcounter{enumi}{2}
\tightlist
\item
  设 \(\psi\) 为双代数 \(\mathsf{TS}(\mathfrak{g})\) 到双代数 \(\mathsf{S}(\mathfrak{g})\) 上的典范同构（《代数》，第 IV 章，第 5 节，命题 12 的推论 1）。映射
\end{enumerate}

\[
\eta \circ \psi : T S (\mathfrak {g}) \rightarrow U
\]

称为典范的。它是唯一的分次余代数态射 \(\eta'\)（TS(g) 到 U 的），使得 \(\eta'(1) = 1\) 且 \(\eta'(x) = \sigma(x)\)（对一切 \(x \in \mathfrak{g}\)）。

\begin{enumerate}
\def\labelenumi{(\arabic{enumi})}
\setcounter{enumi}{3}
\tightlist
\item
  设 V 为向量空间。双代数 \(S(V)\) 的本原元素是次数 1 的元素。这由命题 9 的推论应用于交换李代数 V 得出。
\end{enumerate}

设 \((e_i)_{i\in I}\) 为向量 K-空间 \(g\) 的基，其中指标集 I 是全序的。对一切 \(\alpha \in \mathbf{N}^{(I)}\)，我们记

\[
e _ {\alpha} = \prod_ {i \in I} \frac {\sigma (e _ {i}) ^ {\alpha (i)}}{\alpha (i) !}.\tag{13}
\]

诸 \(e_{\alpha}\)（对 \(|\alpha| \leqslant n\)）构成向量 K-空间 \(\mathbf{U}_n\) 的基（第 I 章，第 2 节，第 7 号，定理 1 的推论 3）。于是

\[
e _ {0} = 1, \quad e _ {\varepsilon_ {i}} = \sigma (e _ {i}) \quad \text { for } i \in I.
\]

由于与滤过代数 U 相伴的分次代数是交换的（同处，定理 1），对 \(\alpha, \beta\) ∈ \(\mathbf{N}^{(I)}\)，

\[
e _ {\alpha}. e _ {\beta} \equiv ((\alpha , \beta)). e _ {\alpha + \beta} \mod U _ {| \alpha | + | \beta | - 1}.\tag{14}
\]

其中

\[
((\alpha , \beta)) = \prod_ {i \in I} \frac {(\alpha (i) + \beta (i)) !}{\alpha (i) ! \beta (i) !}.
\]

另一方面，我们立即有

\[
\varepsilon (e _ {0}) = 1, \quad \varepsilon (e _ {\alpha}) = 0 \quad \text { for } | \alpha | \geqslant 1.\tag{15}
\]

最后，公式 (12) 蕴含：对 \(\alpha \in \mathbf{N}^{(I)}\)，

\[
c (e _ {\alpha}) = \sum_ {\beta + \gamma = \alpha} e _ {\beta} \otimes e _ {\gamma}.\tag{16}
\]

这一公式使我们能确定与余代数 U 对偶的代数 \(\mathrm{U}' = \mathrm{Hom}(\mathrm{U},\mathrm{K})\)（《代数》，第 III 章，第 11 节，第 2 号）。设 \(\mathrm{K}[[\mathrm{X}_i]_{i\in \mathbf{I}}]\) 为不定元 \((\mathrm{X}_i)_{i\in \mathbf{I}}\) 的形式幂级数代数（参看~《代数》，第 III 章，第 2 节，第 11 号）；若 \(\lambda \in \mathrm{U}'\)，用 \(f_{\lambda}\) 表示形式幂级数

\[
f _ {\lambda} = \sum_ {\alpha} \langle \lambda , e _ {\alpha} \rangle \mathrm{X} ^ {\alpha}, \text { where } \mathrm{X} ^ {\alpha} = \prod_ {i \in I} \mathrm{X} _ {i} ^ {\alpha (i)}
\]

其中求和指标 \(\alpha\) 取遍 \(\mathbf{N}^{(I)}\)。

命题 10. 映射 \(\lambda \mapsto f_{\lambda}\) 是代数 \(\mathbf{U}'\) 到形式幂级数代数 \(\mathbf{K}[[\mathbf{X}_i]]_{i \in I}\) 上的同构。

由于 \((e_{\alpha})\) 是 U 的基，映射 \(\lambda \mapsto f_{\lambda}\) 是 K-线性的且双射。另一方面，对 U' 中的 \(\lambda, \mu\)，

\[
\begin{array}{r l} f _ {\lambda \mu} & = \sum_ {\alpha} \langle \lambda \mu , e _ {\alpha} \rangle \mathbf {X} ^ {\alpha} = \sum_ {\alpha} \langle \lambda \otimes \mu , c (e _ {\alpha}) \rangle \mathbf {X} ^ {\alpha} \\ & = \sum_ {\alpha} \langle \lambda \otimes \mu , \sum_ {\beta + \gamma = \alpha} e _ {\beta} \otimes e _ {\gamma} \rangle \mathbf {X} ^ {\alpha} \\ & = \sum_ {\beta , \gamma} \langle \lambda , e _ {\beta} \rangle \langle \mu , e _ {\gamma} \rangle \mathbf {X} ^ {\beta + \gamma} = f _ {\lambda} f _ {\mu}, \end{array}\tag{by (16)}
\]

这表明 \(\lambda \mapsto f_{\lambda}\) 是代数同构，从而完成证明。

\subsection{特征 0 时滤过双代数的结构}

本号中我们继续假设 K 是特征 0 的域。

若 E 是双代数，典范单射 \(P(E) \rightarrow E\) 可以扩张为双代数态射 \(f_{E}: U(P(E)) \rightarrow E\)（第 4 号，命题 3）。

定理 1. 设 E 为余交换双代数。

\begin{enumerate}
\def\labelenumi{(\alph{enumi})}
\item
  双代数态射 \(f_{\mathbf{E}} \colon \mathrm{U}(\mathrm{P}(\mathbf{E})) \to \mathbf{E}\) 是单射。
\item
  若 E 上存在与其双代数结构相容的滤过（第 3 号，定义 2），则态射 \(f_{\mathrm{E}}\) 是同构。
\end{enumerate}

（在情形 (b)，双代数 E 因此等同于其本原元素李代数的包络双代数。）

设 \(c_{\mathbb{E}}\)（相应地，\(\varepsilon_{\mathbb{E}}\)）为 E 的余积（相应地，余单位）。我们记 \(\mathfrak{g} = \mathrm{P}(\mathrm{E})\)；设 \((e_i)_{i \in \mathrm{I}}\) 为向量 K-空间 \(\mathfrak{g}\) 的基，其中指标集 I 是全序的，并设 \((e_\alpha)_{\alpha \in \mathbf{N}^{(I)}}\) 为前一号中引入的基。我们记 \(\mathbf{X}_\alpha = f_{\mathbb{E}}(e_\alpha)\)（对 \(\alpha \in \mathbf{N}^{(I)}\)）。由 (15) 与 (16)，我们有：

\begin{enumerate}
\def\labelenumi{(\arabic{enumi})}
\setcounter{enumi}{16}
\tightlist
\item
\end{enumerate}

\[
\varepsilon_ {E} (X _ {0}) = 1, \quad \varepsilon_ {E} (X _ {\alpha}) = 0 \quad \text { for } | \alpha | \geqslant 1,\tag{18}
\]

\[
c _ {E} (X _ {\alpha}) = \sum_ {\beta + \gamma = \alpha} X _ {\beta} \otimes X _ {\gamma} \quad \text { for } \alpha \in \mathbf {N} ^ {(I)},
\]

因为 \(f_{E}\) 是余代数态射。

我们证明 \(f_{E}\) 是单射。这由下述引理得出：

引理 2. 设 V 为向量空间，E 为余代数，\(f\colon\mathsf{S}(\mathrm{V})\to\mathrm{E}\) 为余代数态射。若 \(f\) 在 \(\mathsf{S}^{\circ}(\mathrm{V})+\mathsf{S}^{1}(\mathrm{V})\) 上的限制是单射，则 \(f\) 是单射。

设 \(n \geqslant 0\)；我们记 \(\mathbb{S}_n = \sum_{i \geqslant n} \mathbb{S}^i(\mathbf{V})\)，用 \(c_s\) 表示 \(\mathbb{S}(\mathbf{V})\) 的余积，并对 \(n\) 用归纳法证明 \(f \mid \mathbb{S}_n\) 是单射。由于断言对 \(n = 0\) 与 \(n = 1\) 是平凡的，我们设 \(n \geqslant 2\)，并设 \(u \in \mathbb{S}_n\) 满足 \(f(u) = 0\)。那么

\[
\begin{array}{r l} {0 = c _ {\mathbf {E}} (f (u))} & {= (f \otimes f) (c _ {\mathbf {s}} (u))} \\ & {= f (u) \otimes 1 + 1 \otimes f (u) + (f \otimes f) (c _ {\mathbf {s}} ^ {+} (u))} \\ & {= (f \otimes f) (c _ {\mathbf {s}} ^ {+} (u)).} \end{array}
\]

由于 \(c_{\mathbb{S}}^{+}(u) \in \mathbb{S}_{n-1} \otimes \mathbb{S}_{n-1}\)，由 (11) 与归纳假设可知 \(u\) 是 \(\mathbb{S}(\mathrm{V})\) 的本原元素，因而是次数 1 的（第 5 号，注 4），从而由于 \(f \mid \mathbb{S}^1(\mathrm{V})\) 是单射而为零。

特别地由此得出族 \((X_{\alpha})\) 是自由的。

我们证明：\(f_{\mathbb{E}}\) 当 \(\mathrm{E}\) 带有与其双代数结构相容的滤过时是满射。设 \((\mathrm{E}_n)_{n\geqslant 0}\) 为这样一个滤过，记 \(\mathrm{E}_n^+ = \mathrm{E}_n\cap \mathrm{Ker}(\epsilon_{\mathbb{E}})\)。我们对 \(n\) 用归纳法证明 \(\mathrm{E}_n^+\) 包含在 \(f_{\mathbb{E}}\) 的像中。由于 \(\mathrm{E} = \mathrm{K}.1 + \bigcup_{n > 0}\mathrm{E}_n^+\)，这将蕴含 \(f_{\mathbb{E}}\) 的满射性。断言对 \(n = 0\) 是平凡的，对 \(n = 1\) 由第 3 号命题 6 的推论得出；以下设 \(n\geqslant 2\)，并设 \(x\in \mathrm{E}_n^+\)。由第 3 号命题 6，

\[
c _ {\mathrm{E}} ^ {+} (x) \in \sum_ {i = 1} ^ {n - 1} \mathrm{E} _ {i} ^ {+} \otimes \mathrm{E} _ {n - i} ^ {+}
\]

且由归纳假设，存在标量 \(\lambda_{\alpha, \beta}\)，其中 \(\alpha, \beta\) 属于 \(\mathbf{N}^{(I)}\)，除有限多个外均为零，使得

\[
c _ {\mathbb {E}} ^ {+} (x) = \sum_ {\alpha , \beta \neq 0} \lambda_ {\alpha , \beta} \mathrm{X} _ {\alpha} \otimes \mathrm{X} _ {\beta}.\tag{19}
\]

于是由公式 (18)

\[
(c _ {\mathrm{E}} ^ {+} \otimes \operatorname{Id} _ {\mathrm{E}}) (c _ {\mathrm{E}} ^ {+} (x)) = \sum_ {\alpha , \beta , \gamma \neq 0} \lambda_ {\alpha + \beta , \gamma} \mathbf {X} _ {\alpha} \otimes \mathbf {X} _ {\beta} \otimes \mathbf {X} _ {\gamma}
\]

\[
(\mathrm{Id} _ {\mathtt {E}} \otimes c _ {\mathtt {E}} ^ {+}) (c _ {\mathtt {E}} ^ {+} (x)) = \sum_ {\alpha , \beta , \gamma \neq 0} \lambda_ {\alpha , \beta + \gamma} \mathbf {X} _ {\alpha} \otimes \mathbf {X} _ {\beta} \otimes \mathbf {X} _ {\gamma}.
\]

由第 1 号命题 3 以及诸 \(X_{\alpha}\) 的线性无关性，得出

\[
\lambda_ {\alpha + \beta , \gamma} = \lambda_ {\alpha , \beta + \gamma} \quad \text { for } \alpha , \beta , \gamma \text { in } \mathbf {N} ^ {(I)} - \{0 \}.\tag{20}
\]

另一方面，余积 \(c_{E}\) 是余交换的；与上面同样的论证蕴含

\[
\lambda_ {\alpha , \beta} = \lambda_ {\beta , \alpha} \quad \text { for } \alpha , \beta \text { in } \mathbf {N} ^ {(I)} - \{0 \}.\tag{21}
\]

假设存在标量族 \((\mu_{\alpha})\)（\(|\alpha| \geqslant 2\)），使得

\[
\mu_ {\alpha + \beta} = \lambda_ {\alpha , \beta} \quad \text { for } \alpha , \beta \text { in } \mathbf {N} ^ {(I)} - \{0 \}.\tag{22}
\]

那么

\[
c _ {\mathrm{E}} ^ {+} (x) = \sum_ {\alpha , \beta \neq 0} \mu_ {\alpha + \beta} \mathrm{X} _ {\alpha} \otimes \mathrm{X} _ {\beta} = \sum_ {| \gamma | \geqslant 2} \mu_ {\gamma} c _ {\mathrm{E}} ^ {+} (\mathrm{X} _ {\gamma})
\]

由公式 (18)，于是 \(y = x - \sum_{|\gamma| \geq 2} \mu_\gamma X_\gamma\) 是本原的，因而属于 \(P(E) \subset \operatorname{Im}(f_E)\)。从而

\[
x = y + \sum_ {| \gamma | \geqslant 2} \mu_ {\gamma} f _ {\mathrm{E}} (e _ {\gamma}) \in \operatorname{Im} (f _ {\mathrm{E}}).
\]

因此，只要证明了下述引理，证明即告完成：

引理 3. 若有限支集的标量族 \((\lambda_{\alpha, \beta})\)（其中 \(\alpha, \beta\) 属于 \(\mathbf{N}^{(I)} - \{0\}\)）满足关系 (20) 与 (21)，则存在有限支集的族 \((\mu_{\alpha})_{|\alpha| \geqslant 2}\)，使得 \(\mu_{\alpha + \beta} = \lambda_{\alpha, \beta}\) 对非零的 \(\alpha, \beta\) 成立。

只需证明

\[
\alpha + \beta = \gamma + \delta\tag{23}
\]

蕴含 \(\lambda_{\alpha, \beta} = \lambda_{\gamma, \delta}\) 对非零的 \(\alpha, \beta, \gamma, \delta\) 成立。由 Riesz 分解引理（《代数》，第 VI 章，第 1 节，第 10 号，定理 1），存在 \(\pi, \rho, \sigma, \tau\)（\(\mathbf{N}^{(I)}\) 中的），使得

\[
\alpha = \pi + \sigma , \quad \beta = \rho + \tau , \quad \gamma = \pi + \rho , \quad \delta = \sigma + \tau .
\]

设 \(\pi \neq 0\)；由于 \(\sigma + \beta = \rho + \delta\)，关系 (20) 蕴含

\[
\lambda_ {\alpha , \beta} = \lambda_ {\pi + \sigma , \beta} = \lambda_ {\pi , \sigma + \delta} = \lambda_ {\pi , \rho + \delta} = \lambda_ {\pi + \rho , \delta} = \lambda_ {\gamma , \delta}.
\]

反之，若 \(\pi = 0\)，则 \(\beta = \gamma + \tau\) 且 \(\delta = \alpha + \tau\)，由此

\[
\lambda_ {\alpha , \beta} = \lambda_ {\alpha , \gamma + \tau} = \lambda_ {\alpha + \tau , \gamma} = \lambda_ {\delta , \gamma}
\]

由 (20) 得出，而又由 (21) 有 \(\lambda_{\delta, \gamma} = \lambda_{\gamma, \delta}\)，故 \(\lambda_{\alpha, \beta} = \lambda_{\gamma, \delta}\)。

\section{自由李代数}

\subsection{自由代数的复习}

设 X 为集合。回忆在 X 上构造的自由原群 M(X) 的作法（《代数》，第 I 章，第 7 节，第 1 号）。对整数 \(n \geqslant 1\) 用归纳法，我们如下定义诸集合 \(\mathbf{X}_n\)：写 \(\mathbf{X}_1 = \mathbf{X}\)，并取 \(\mathbf{X}_n\) 为诸集合 \(\mathbf{X}_p \times \mathbf{X}_{n-p}\)（其中 \(p = 1, 2, \ldots, n-1\)）的和集；若 X 有限，则每个 \(\mathbf{X}_n\) 也有限。族 \((\mathbf{X}_n)_{n \geqslant 1}\) 的和集记作 M(X)；每个集合 \(\mathbf{X}_n\)（特别地 X）都等同于 M(X) 的一个子集。 设 w 与 \(w'\) 属于 M(X)；设 p 与 q 为使 \(w \in \mathbf{X}_p\) 与 \(w' \in \mathbf{X}_q\) 的整数，并设 \(n = p + q\)；有序对 \(\langle w, w' \rangle\) 在 \(\mathbf{X}_p \times \mathbf{X}_{n-p}\) 到 \(\mathbf{X}_n\) 的典范单射下的像记作 \(w.w'\)，称为 w 与 \(w'\) 的乘积。 X 到原群 M 的每个映射都能以唯一方式扩张为 M(X) 到 M 的原群同态。

设 \(w\) 属于 \(\mathrm{M}(\mathbf{X})\)；唯一整数 \(n\)（使 \(w \in \mathbf{X}_n\) 的）称为 \(w\) 的长度，记作 \(l(w)\)。于是 \(l(w.w') = l(w) + l(w')\) 对 \(w, w'\) ∈ \(\mathrm{M}(\mathbf{X})\) 成立。 集合 \(\mathbf{X}\) 是 \(\mathrm{M}(\mathbf{X})\) 中由长度为 1 的元素组成的子集。每个元素 \(w\)（长度 \(\geqslant 2\) 的）都能唯一地写成 \(w = w'.w''\) 的形式。

环 K 上系数的原群 M(X) 的代数记作 \(\operatorname{Lib}(X)\)，当需要指明环 K 时记作 \(\operatorname{Lib}_{K}(X)\)。集合 M(X) 是 K-模 \(\operatorname{Lib}(X)\) 的基，因此 X 将等同于 \(\operatorname{Lib}(X)\) 的一个子集。若 A 是代数，X 到 A 的每个映射都可唯一地扩张为 \(\operatorname{Lib}(X)\) 到 A 的同态（《代数》，第 III 章，第 2 节，第 7 号，命题 7）。

\subsection{自由李代数的构造}

定义 1. 集合 \(\mathbf{X}\) 上的自由李代数是商代数

\[
\mathrm{L} (\mathrm{X}) = \operatorname{Lib} (\mathrm{X}) / \mathrm{a},
\]

其中 a 是 \(\operatorname{Lib}(\mathbf{X})\) 的由具有下列形式之一的元素生成的双边理想

\begin{enumerate}
\def\labelenumi{(\arabic{enumi})}
\tightlist
\item
\end{enumerate}

\[
\mathrm{Q} (a) = a. a \text {   for   } a \text {   in   } \operatorname{Lib} (\mathrm{X}),\tag{2}
\]

\[
J (a, b, c) = a. (b. c) + b. (c. a) + c. (a. b)
\]

其中 \(a, b, c\) 属于 \(\operatorname{Lib}(\mathbf{X})\)。

显然 \(\mathrm{L}(\mathrm{X})\) 是李 K-代数；两个元素 \(u, v\)（\(\mathrm{L}(\mathrm{X})\) 的元素）的合成将记作 \([u, v]\)。当需要指明环 \(\mathbf{K}\) 时，我们用 \(\mathrm{L}_{\mathbb{K}}(\mathbf{X})\) 记 \(\mathrm{L}(\mathbf{X})\)。

下述命题说明了赋予 L(X) 自由李代数这一名称的理由。

命题 1. 设 \(\psi\) 为 \(\operatorname{Lib}(\mathbf{X})\) 到 \(\operatorname{L}(\mathbf{X})\) 上的典范映射，\(\phi\) 为 \(\psi\) 在 \(\mathbf{X}\) 上的限制。对每个映射 \(f\)（\(\mathbf{X}\) 到李代数 \(\mathfrak{g}\) 的映射），存在且仅存在一个同态 \(\mathrm{F}:\mathrm{L}(\mathrm{X})\to \mathfrak{g}\)，使得 \(f = \mathrm{F}\circ \phi\)。

\begin{enumerate}
\def\labelenumi{(\alph{enumi})}
\item
  F 的存在性：设 h 为 \(\operatorname{Lib}(\mathbf{X})\) 到 g 的扩张 \(f\) 的同态（第 1 号）。对 \(\operatorname{Lib}(\mathbf{X})\) 中的一切 a，\(h(\mathbf{Q}(a)) = h(a.a) = [h(a), h(a)] = 0\)； 类似地，g 所满足的 Jacobi 恒等式蕴含 \(h(\mathbf{J}(a, b, c)) = 0\)（对 \(\operatorname{Lib}(\mathbf{X})\) 中的 a, b, c）。 由此得出 \(h(a) = 0\)，从而存在 \(\operatorname{L}(\mathbf{X})\) 到 g 的同态 F，使得 \(h = F \circ \phi\)。 限制到 X 上，即得 \(f = F \circ \phi\)。
\item
  F 的唯一性：设 \(\mathrm{F}^{\prime}:\mathrm{L}(\mathrm{X})\to\mathfrak{g}\) 为使 \(f=\mathrm{F}^{\prime}\circ\phi\) 的同态。同态 \(F\circ\psi\) 与 \(F^{\prime}\circ\psi\)（\(\operatorname{Lib}(X)\) 到 g 的）在 X 上一致，因而相等；由于 \(\psi\) 是满射，\(F=F^{\prime}\)。
\end{enumerate}

推论 1. 族 \((\phi (x))_{x\in \mathbf{X}}\) 在 \(\mathbf{K}\) 上是 \(\mathrm{L}(\mathbf{X})\) 中的自由族。

设 \(x_{1}, x_{2}, \ldots, x_{n}\) 为 X 中互异的元素，\(\lambda_{1}, \ldots, \lambda_{n}\) 为 K 中满足下式的元素

\[
\lambda_ {1}. \phi (x _ {1}) + \dots + \lambda_ {n}. \phi (x _ {n}) = 0.\tag{3}
\]

设 \(\mathfrak{g}\) 为以 \(\mathbf{K}\) 为基础模的交换李代数。对 \(i = 1,2,\ldots ,n\)，存在同态 \(\mathrm{F}_i\)（\(\mathbf{L}(\mathbf{X})\) 到 \(\mathfrak{g}\) 的），使得 \(\mathrm{F_i}(\phi (x_i)) = 1\) 且 \(\mathrm{F_i}(\phi (x)) = 0\) 对 \(x\neq x_{i}\) 成立（命题 1）；把 \(\mathrm{F}_i\) 应用于关系 (3)，便得 \(\lambda_{i} = 0\)。

推论 2. 设 \(\mathfrak{a}\) 为李代数。\(\mathrm{L}(\mathrm{X})\) 被 \(\mathfrak{a}\) 的每个扩张都是非本质的。

设 \(\mathfrak{a} \xrightarrow{\lambda} \mathfrak{g} \xrightarrow{\mu} \mathrm{L}(\mathrm{X})\) 为这样一个扩张（第 I 章，第 1 节，第 7 号）。由于 \(\mu\) 是满射，存在映射 \(f\)（\(\mathrm{X}\) 到 \(\mathfrak{g}\) 的），使得 \(\phi = \mu \circ f\)。 设 \(\mathrm{F}\) 为 \(\mathrm{L}(\mathrm{X})\) 到 \(\mathfrak{g}\) 的使 \(f = \mathrm{F} \circ \phi\) 的同态（命题 1）。 于是 \((\mu \circ \mathrm{F}) \circ \phi = \mu \circ f = \phi\)，命题 1 表明 \(\mu \circ \mathrm{F}\) 是 \(\mathrm{L}(\mathrm{X})\) 的恒同自同构。 因此所给的扩张是非本质的（第 I 章，第 1 节，第 7 号，命题 6 与定义 6）。

由于环 K 非零，命题 1 的推论 1 表明 \(\phi\) 是单射。因此集合 X 可借助 \(\phi\) 等同于它在 L(X) 中的像；在这一约定下，X 生成 L(X)，且 X 到李代数 \(\mathfrak{g}\) 的每个映射都可扩张为 L(X) 到 \(\mathfrak{g}\) 的李代数同态。

注. 当 X 为空时，M(X) 为空，因而 L(X) = \{0\}。若 X 由单个元素 x 组成，则 L(X) 的子模 K.x 是

\(L(X)\) 的李子代数；由于 X 生成 \(L(X)\)，命题 1 的推论 1 表明 \(L(X)\) 是以 \(\{x\}\) 为基的自由模。

\subsection{李代数的表现}

设 \(\mathfrak{g}\) 为李代数，\(\pmb{a} = (a_i)_{i\in I}\) 为 \(\mathfrak{g}\) 的元素族。设 \(f_{a}\) 为 \(\mathrm{L(I)}\) 到 \(\mathfrak{g}\) 的把每个 \(i\in I\) 映为 \(a_{i}\) 的同态。\(f_{a}\) 的像是 \(\mathfrak{g}\) 中由 \(\pmb{a}\) 生成的子代数；\(f_{a}\) 的核的元素称为族 \(\pmb{a}\) 的关系元。族 \(\pmb{a}\) 若 \(f_{a}\) 是满射（相应地，单射、双射），则称为生成的（相应地，自由的、基的）。

设 \(\mathfrak{g}\) 为李代数。\(\mathfrak{g}\) 的一个表现是指一个有序对 \((\boldsymbol{a},\boldsymbol{r})\)，由生成族 \(\boldsymbol{a} = (a_{i})_{i\in \mathbb{I}}\) 与关系元族 \(\boldsymbol{r} = (r_j)_{j\in \mathbb{J}}\)（\(\boldsymbol{a}\) 的关系元的族）组成，后者生成 L(I) 的理想即 \(f_{a}\) 的核。我们也说 \(\mathfrak{g}\) 由族 \(\boldsymbol{a}\) 连同关系元 \(r_j\)（\(j\in J\)）所表现。

设 I 为集合，\(\pmb{r} = (r_j)_{j \in J}\) 为自由李代数 L(I) 的元素族；设 \(a_r\) 为 L(I) 的由 \(\pmb{r}\) 生成的理想。商代数 L(I, \(\pmb{r}\) ) = L(I)/ \(a_r\) 称为由 I 与关系元族（\(r_j)_{j \in J}\)）所定义的李代数；我们也说 L(I, \(\pmb{r}\) ) 由表现 (I, \(\pmb{r}\) ) 定义，或由 (I; ( \(r_j = 0\) ) \(_{j \in J}\)) 定义。当族 \(\pmb{r}\) 为空时，L(I, \(\pmb{r}\) ) = L(I)。

设 \(I\) 与 \(r\) 如上；设 \(\xi_i\) 表示 \(i\) 在 \(L(I, r)\) 中的像。生成族 \(\xi = (\xi_i)_{i \in I}\) 与关系元族 \(r\) 构成 \(L(I, r)\) 的一个表现。反之，若 \(g\) 为李代数且 \((a, r)\)（其中 \(a = (a_i)_{i \in I}\)）是 \(g\) 的一个表现，则存在唯一的同构 \(u: L(I, r) \to g\)，使得 \(u(\xi_i) = a_i\) 对一切 \(i \in I\) 成立。

\subsection{李多项式与代入}

设 I 为集合。设 \(\mathrm{T}_i\) 表示 I 的元素 \(i\) 在 \(\mathbf{L}(\mathbf{I})\) 中的典范像（后者有时也记作 \(\mathrm{L}((\mathrm{T}_i)_{i\in \mathbf{I}})\)）；\(\mathrm{L}(\mathbf{I})\) 的元素称为不定元 \((\mathrm{T}_i)_{i\in \mathbf{I}}\) 的李多项式。

设 g 为李代数。若 \(\boldsymbol{t} = (t_{i})_{i \in I}\) 是 g 的元素族，用 \(f_{\boldsymbol{t}}\) 表示 L(I) 到 g 的同态，使得 \(f_{\boldsymbol{t}}(\mathrm{T}_{\boldsymbol{i}}) = t_{i}\)（对 \(i \in I\)）（第 2 号，命题 1）。L(I) 的元素 P 在 \(f_{\boldsymbol{t}}\) 下的像记作 \(\mathrm{P}((t_{i})_{i \in I})\)。特别地，\(\mathrm{P}((\mathrm{T}_{\boldsymbol{i}})_{i \in I}) = \mathrm{P}\)；上述元素 \(\mathrm{P}((t_{i})_{i \in I})\) 有时称为以 \(t_{i}\) 代入 \(T_{i}\) 到李多项式 \(\mathrm{P}((\mathrm{T}_{\boldsymbol{i}})_{i \in I})\) 中所得的 g 的元素。

设 \(\sigma: \mathfrak{g} \to \mathfrak{g}'\) 为李代数同态。对每个族 \(t = (t_i)_{i \in I}\)（\(\mathfrak{g}\) 的元素的族）以及所有 \(P \in L(I)\)，

\[
\sigma (\mathrm{P} ((t _ {i}) _ {i \in \mathrm{I}})) = \mathrm{P} ((\sigma (t _ {i})) _ {i \in \mathrm{I}}),\tag{4}
\]

因为 \(\sigma^{\circ}f_{t}\) 把 \(\mathrm{T}_i\) 映为 \(\sigma(t_i)\)（对 \(i\in I\)）。

设 \((\mathbf{Q}_j)_{j\in J}\) 为 \(\mathbf{L}(\mathbf{I})\) 的元素族，并设 \(\mathbf{P}\in \mathbf{L}(\mathbf{J})\)。以 \(Q_{j}\) 代入 \(T_{j}\)（在 \(P\) 中），得到李多项式 \(R = P((Q_{j})_{j \in J}) \in L(I)\)。那么

\[
\mathrm{R} ((t _ {i}) _ {i \in \mathrm{I}}) = \mathrm{P} ((Q _ {j} ((t _ {i}) _ {i \in \mathrm{I}})) _ {j \in J})\tag{5}
\]

对每个族 \(\pmb{t} = (t_i)_{i\in \mathbf{I}}\)（李代数 \(\mathfrak{g}\) 的元素的族）成立，这可由用同态 \(f_{t}\) 作用于等式 \(\mathrm{R} = \mathrm{P}((\mathrm{Q}_j)_{j\in \mathbf{J}})\) 并利用 (4) 看出。

设 \(\mathfrak{g}\) 为李代数，I 为有限集，\(P \in L(I)\)。假设 \(\mathfrak{g}\) 是自由 K-模。映射

\[
\tilde {\mathrm{P}} \colon \mathfrak {g} ^ {\mathrm{I}} \to \mathfrak {g}
\]

由 \(\tilde{\mathrm{P}}((t_i)_{i\in \mathbf{I}}) = \mathrm{P}((t_i)_{i\in \mathbf{I}})\) 定义的映射于是是一个多项式映射。†因为 \(\mathfrak{g}^{\mathbf{I}}\) 到 \(\mathfrak{g}\) 的映射的集合 F 按由下式定义的括号成一个李代数

\[
[ \phi , \psi ] (\boldsymbol {t}) = [ \phi (\boldsymbol {t}), \psi (\boldsymbol {t}) ];\tag{6}
\]

多项式映射的集合 \(\mathbf{F}'\)（\(\mathfrak{g}^{\intercal}\) 到 \(\mathfrak{g}\) 的） 由括号的双线性是 \(\mathbf{F}\) 的李子代数。于是我们的断言由如下事实得出：映射 \(\mathrm{P} \mapsto \tilde{\mathrm{P}}\) 是李代数同态，且 \(\tilde{\mathrm{T}}_i = \mathrm{pr}_i \in \mathrm{F}'\) 对一切 \(i\) 成立。

\subsection{函子性质}

命题 2. 设 X 与 Y 为两个集合。每个映射 \(u: \mathrm{X} \to \mathrm{Y}\) 都可唯一地扩张为李代数同态 \(\mathrm{L}(u): \mathrm{L}(\mathrm{X}) \to \mathrm{L}(\mathrm{Y})\)。对每个映射 \(v: \mathrm{Y} \to \mathrm{Z}\)，有 \(\mathrm{L}(v \circ u) = \mathrm{L}(v) \circ \mathrm{L}(u)\)。

\(\mathrm{L}(u)\) 的存在性与唯一性由第 2 号命题 1 得出。同态 \(\mathrm{L}(v\circ u)\) 与 \(\mathrm{L}(v)\circ \mathrm{L}(u)\) 在 X 上有相同的限制，因而相等（命题 1）。

推论. 若 \(u\) 是单射（相应地，满射、双射），则 \(L(u)\) 亦然。

由于断言对 \(\mathrm{X} = \varnothing\) 是平凡的，我们设 \(\mathrm{X} \neq \varnothing\)。若 \(u\) 是单射，则存在映射 \(v\)（\(\mathrm{Y}\) 到 \(\mathrm{X}\) 的），使 \(v \circ u\) 是 \(\mathrm{X}\) 的恒同映射；由命题 2，\(\mathrm{L}(v) \circ \mathrm{L}(u)\) 是 \(\mathrm{L}(\mathrm{X})\) 的恒同自同构，因而 \(\mathrm{L}(u)\) 是单射。当 \(u\) 是满射时，存在映射 \(w\)（\(\mathrm{Y}\) 到 \(\mathrm{X}\) 的），使 \(u \circ w\) 是 \(\mathrm{Y}\) 的恒同映射；于是 \(\mathrm{L}(u) \circ \mathrm{L}(w)\) 是 \(\mathrm{L}(\mathrm{Y})\) 的恒同映射，这就证明 \(\mathrm{L}(u)\) 是满射。

设 \(\mathbf{X}\) 为集合，\(\mathbf{S}\) 为 \(\mathbf{X}\) 的子集。上述推论表明，\(\mathbf{S}\) 到 \(\mathbf{X}\) 的典范单射可扩张为同构 \(\alpha\)（\(\mathbf{L}(\mathbf{S})\) 到李子代数 \(\mathbf{L}'(\mathbf{S})\) 上的同构，后者是 \(\mathbf{L}(\mathbf{X})\) 中由 \(\mathbf{S}\) 生成的李子代数）；我们将把 \(\mathbf{L}(\mathbf{S})\) 与 \(\mathbf{L}'(\mathbf{S})\) 借助 \(\alpha\) 等同起来。

设 \((\mathrm{S}_{\alpha})_{\alpha \in \mathbf{I}}\) 为 X 的子集的一个右向族，其并为 S。关系 \(\mathrm{S}_{\alpha} \subset \mathrm{S}_{\beta}\) 蕴含 \(\mathrm{L}(\mathrm{S}_{\alpha}) \subset \mathrm{L}(\mathrm{S}_{\beta})\)，因而李子代数族 \(\mathrm{L}(\mathrm{S}_{\alpha})\)（\(\mathrm{L}(\mathrm{X})\) 的）是右向的。因此 \(\mathfrak{g} = \bigcup_{\alpha \in \mathbf{I}} \mathrm{L}(\mathrm{S}_{\alpha})\) 是 \(\mathrm{L}(\mathrm{X})\) 的李子代数；于是 \(\mathrm{S} \subset \mathfrak{g}\)，由此 \(\mathrm{L}(\mathrm{S}) \subset \mathfrak{g}\)，又由于 \(\mathrm{L}(\mathrm{S}_{\alpha}) \subset \mathrm{L}(\mathrm{S})\) 对一切 \(\alpha \in \mathbf{I}\)，故 \(\mathfrak{g} \subset \mathrm{L}(\mathrm{S})\)。于是

\[
\mathrm{L} \left(\bigcup_ {\alpha \in \mathrm{I}} \mathrm{S} _ {\alpha}\right) = \bigcup_ {\alpha \in \mathrm{I}} \mathrm{L} (\mathrm{S} _ {\alpha})\tag{7}
\]

对 X 的子集的每个右向族 \((\mathrm{S}_{\alpha})_{\alpha\in\mathbb{I}}\) 成立。

把上述结果应用于 X 的有限子集的族，我们看到 \(\mathrm{L}(\mathrm{X})\) 的每个元素都具有形式 \(\mathrm{P}(x_{1},\ldots,x_{n})\)，其中 P 是 n 个不定元的李多项式，而 \(x_{1},\ldots,x_{n}\) 是 X 的元素。

命题 3. 设 \(\mathbf{K}'\) 为非零交换环，\(u: \mathbf{K} \to \mathbf{K}'\) 为环同态。对每个集合 \(\mathbf{X}\)，存在且仅存在一个李 \(\mathbf{K}'\) -代数同态

\[
v \colon \mathrm{L} _ {\mathbf {K}} (\mathrm{X}) \otimes \mathrm{K} ^ {\prime} \rightarrow \mathrm{L} _ {\mathbf {K} ^ {\prime}} (\mathrm{X})
\]

使得 \(v(x \otimes 1) = x\) 对 \(x \in \mathbf{X}\) 成立。而且，\(v\) 是同构。

把命题 1 应用于视为李 K-代数的 \(\mathfrak{g} = \mathrm{L}_{\mathbb{K}'}(\mathrm{X})\) 以及映射 \(x\mapsto x\)（\(\mathbf{X}\) 到 \(\mathfrak{g}\) 的），我们得到一个 K-同态 \(\mathrm{L}_{\mathbb{K}}(\mathrm{X})\to \mathrm{L}_{\mathbb{K}'}(\mathrm{X})\)，从而存在 K'-同态 \(v\colon \mathrm{L}_{\mathbb{K}}(\mathrm{X})\otimes \mathrm{K}'\to \mathrm{L}_{\mathbb{K}'}(\mathrm{X})\)。\(v\) 的唯一性以及它是同构这一事实，由有序对 \((\mathrm{L}_{\mathbb{K}}(\mathrm{X})\otimes \mathrm{K}',x\mapsto x\otimes 1)\) 与有序对 \((\mathrm{L}_{\mathbb{K}'}(\mathrm{X}),x\mapsto x)\) 是同一泛映射问题的解这一事实得出。

注. 设 \(\mathfrak{h}'\) 为李 K'-代数，\(\mathfrak{h}\) 为由 \(\mathfrak{h}'\) 通过限制标量环而得到的李 K-代数。若 \(\mathrm{P} \in \mathrm{L}_{\mathrm{K}}(\mathrm{X})\)，我们可以定义 \(\tilde{\mathrm{P}}: \mathfrak{h}^{\mathrm{X}} \to \mathfrak{h}\)（第 4 号）。我们立即看出

\[
\tilde {\mathrm{P}} = (v (\mathrm{P} \otimes 1)) ^ {\sim}.
\]
